\chapter{El amparo}

\label{cap:amparo}
El 21 de diciembre de 2018 y el 3 de enero de 2019, dos aluviones sobre la cuenca del río La Caldera dañaron
viviendas de loteos construidos junto a sus arroyos. De ahí salió un amparo colectivo que lleva más de siete años y medio en trámite, en una vía que la Constitución define como expedita, y que produjo la pieza documental más importante de este libro: una sentencia
que declara la obligación estatal de proteger y recomponer la cuenca, con criterio
ecocéntrico y cita de la Corte Interamericana.

Este capítulo sigue esa causa de principio a fin: qué ordenó, cómo se ejecutó, qué
dijeron las partes al defenderse y qué pasó en el mes en que el plan finalmente se
aprobó. El capítulo anterior ya estableció qué sabía el Estado sobre esos arroyos antes
de que se desbordaran.

\section{El expediente}

\begin{ficha}{Causa Lazarte Vigabriel --- Amparo Colectivo Ambiental}
\begin{itemize}[leftmargin=*,nosep]
\item \textbf{Autos:} \textit{Lazarte Vigabriel, Eduardo José Ignacio y otros c/ Municipalidad de La Caldera; Ministerio de la Producción, Trabajo y Desarrollo Sustentable de la Provincia de Salta --- Amparo Colectivo}. Expte.\ Nº EMI 23860/2020; la etapa de ejecución tramita como Expte.\ Nº 780280.
\item \textbf{Tribunal:} Juzgado de Minas y en lo Comercial de Registro del Distrito Judicial Centro, a cargo de la jueza \textbf{María Victoria Mosmann}.
\item \textbf{Actor:} Eduardo José Ignacio Lazarte Vigabriel, en condición de afectado, acompañando listado de damnificados y justificando la representación adecuada en su experiencia previa en amparos ambientales colectivos.
\item \textbf{Hechos:} inundaciones de diciembre de 2018 y enero de 2019 en los loteos \textbf{El Durazno I a V, El Nogalar I, II y IV, y Potreros de La Caldera}, por aluvión proveniente de la sedimentación de los cauces de los arroyos \textbf{Guaranguay, Los Duraznos y Los Manzanos}, que integran la microcuenca del río La Caldera. La causa se atribuye a las condiciones geomorfológicas de la cuenca, a las actividades antrópicas allí desplegadas y a la omisión de las autoridades en tomar medidas de prevención y contención.
\item \textbf{Vía:} amparo ambiental colectivo del artículo 43 de la Constitución Nacional, con legitimación fundada en la Ley provincial 7070 y en la Ley General del Ambiente 25.675.
\item \textbf{Fundamento de la obligación:} artículo 41 de la Constitución Nacional, Constitución provincial y leyes provinciales \textbf{7070, 7017 y 7543}.
\end{itemize}
\textit{Fuente: ficha del caso en el Observatorio Argentino de Litigación Climática, Facultad de Ciencias Jurídicas y Sociales de la Universidad Nacional del Litoral. Análisis de Jimena Bevilacqua, Clarisa Neuman y Lorena Bianchi; ficha elaborada por Hernán Caravario, 5 de marzo de 2023.}
\end{ficha}

\begin{ficha}{Lo que el escrito de febrero de 2019 invocaba}
La presentación, con patrocinio del abogado Lazarte ante el Juzgado de Minas, se apoyaba en \textbf{informes técnicos públicos previos}, que el escrito hace corresponder con otros del \textbf{INTA} y de organizaciones ambientales: la \textbf{Asociación Bosque Modelo Jujuy} y la \textbf{Fundación Oikos para el Desarrollo}. Según el escrito, esas instituciones \textbf{ya habían advertido la alta peligrosidad que revestía la cuenca sin manejo apropiado}, dadas sus características de cuenca de montaña con arroyos y ríos de pendiente elevada, sumados a los impactos antrópicos autorizados por la Municipalidad de La Caldera para actividades industriales, mineras y de desarrollo urbanístico.

El área de afectación comprende los caminos vecinales que cruzan los arroyos \textbf{Guaranguay, Los Duraznos, Los Manzanos y uno sin identificación}, todos de paso obligado hacia El Nogalar I, II y IV y hacia \textbf{todas las etapas de El Durazno}, y las actividades en la playa del río sobre la avenida Costanera sur.
\end{ficha}

\begin{cautela}
Esta pieza refuerza lo que el capítulo \ref{cap:pdua} documentó por otra vía y conviene enunciarlo junto: \textbf{la advertencia técnica sobre la cuenca es anterior al daño y no provenía de los vecinos}. Provenía del INTA y de dos organizaciones especializadas, y estaba publicada. El amparo no inaugura el conocimiento del riesgo: lo invoca como preexistente.

Dos precisiones documentales menores, que se dejan asentadas porque el libro no debe fijar cifras que no verificó. Primera: la cobertura de 2019 enumera las etapas de El Durazno como \textbf{I a VI}, mientras la de 2021 y la ficha del Observatorio consignan \textbf{I a V}. La diferencia importa, porque el capítulo cuenta etapas con y sin plano aprobado. Segunda: uno de los cuatro arroyos del área \textbf{carece de identificación} en el propio escrito, lo que es en sí un dato sobre el estado del relevamiento hídrico de la microcuenca.
\end{cautela}

\section{Un caso de litigación climática}
La causa está registrada en el Observatorio Argentino de Litigación Climática, y esa calificación no es decorativa: describe cómo razona la sentencia.

Según la ficha del Observatorio, el tribunal introduce la cuestión del cambio climático \textbf{de modo explícito}, como un hecho notorio que en combinación con la omisión estatal tiene el potencial de generar afectaciones a los derechos humanos. Para ello recurre a la \textbf{Opinión Consultiva OC-23/2017 de la Corte Interamericana de Derechos Humanos} y a la \textbf{Observación General Nº 15 del Comité de Derechos Económicos, Sociales y Culturales}, que menciona los efectos del cambio climático como factor que los Estados deben considerar al establecer sus políticas de acceso al agua. La sentencia emplea además la teoría de los bienes comunes, tratando al ambiente como macrobien y al recurso hídrico de la cuenca como microbien ambiental, y aborda la cuenca como unidad de gestión desde un paradigma ecocéntrico.

\begin{cautela}
Parte de la cobertura periodística consignó que la sentencia funda la legitimación en el Acuerdo de Escazú. La ficha académica del caso no lo menciona: sitúa la legitimación en el artículo 43 de la Constitución Nacional, la ley provincial 7070 y la Ley General del Ambiente, y el andamiaje convencional en la OC-23/2017 y la Observación General Nº 15. Este libro sigue la ficha, que es la fuente más próxima al texto, y deja consignada la divergencia. Verificarla requiere el texto íntegro de la sentencia.
\end{cautela}

\section{Lo que la sentencia ordena}

La sentencia declara la obligación estatal de protección y recomposición de la cuenca y ordena la confección de un \textbf{Plan de Manejo Integral de la microcuenca del río La Caldera}, con base en el principio de progresividad, fijando el contenido de sus etapas y los criterios que los organismos deben seguir.

\textbf{Primera etapa --- línea de base socioambiental o diagnóstico, seis meses.} Determinar las áreas de zona urbana y no urbana comprendidas en la microcuenca; delimitar el área de aplicación mediante cartografía digital que coordine las acciones con el ordenamiento territorial del municipio; definir objetivos generales sobre la normativa ambiental vigente y específicos para la superficie de interés; identificar las externalidades hídricas negativas de elementos naturales o antrópicos que inciden en la fragilidad de la cuenca; identificar y caracterizar las acciones actuales, sus causas y efectos y el manejo de los recursos naturales; y \textbf{considerar una estrategia de estabilización de la microcuenca priorizando la realización de infraestructura verde por sobre infraestructura gris}, con un balance progresivo de recuperación y saneamiento bajo un paradigma ecológico de base constitucional y convencional.

\textbf{Segunda etapa --- el Plan.} Elaborado sobre la información de la primera, con etapas y plazos presentados durante el trámite de cumplimiento. Debe realizarse atendiendo aspectos políticos, físicos, sociales, tecnológicos, culturales, económicos, jurídicos y ecológicos de la realidad local, regional y nacional; asegurar el uso ambientalmente adecuado de los recursos, garantizar la mínima degradación y \textbf{promover la participación social en las decisiones fundamentales del desarrollo sostenible}; y ajustarse a escala predial de la cuenca, al ordenamiento territorial del municipio y a la Ley provincial de Ordenamiento Territorial de Bosques Nativos, recordando que ésta establece como unidad estructural y espacial de análisis la \textbf{cuenca hidrográfica}, el porcentaje de pendiente, y detalla los sectores de muy alto valor de conservación.

El último lineamiento vale una observación. La sentencia manda ajustar el Plan al ordenamiento territorial del municipio: es decir, al Código de Ordenamiento Territorial cuyo artículo 42 clasifica precisamente a los tres loteos actores como núcleos espontáneos, y cuya Zona PU declara no urbanizable la ribera del río.

El tribunal remite al instrumento municipal como marco del Plan, sin que conste en la ficha que se haya pronunciado sobre la brecha entre ese instrumento y su aplicación. El Plan aprobado en 2026 debería, si sigue el lineamiento, resolver esa brecha; contrastarlo contra el articulado del Código es la verificación pendiente más productiva del capítulo.

Y el criterio de pendiente que la sentencia recuerda del OTBN conecta con lo que el capítulo \ref{cap:pdua} documentó: el estudio de cota máxima edificable que el propio plan urbano de 2015 dejó como proyecto pendiente.

Un matiz sobre la fuerza del decisorio. La ficha del Observatorio, al sintetizar, dice que la sentencia ``declara la obligación de protección y recomposición de la cuenca y \textit{exhorta} a los organismos estatales demandados a confeccionar un plan integral'', y más adelante que ``ordena la confección'' del Plan. La diferencia entre exhortar y condenar no es menor en la práctica de ejecución, y explica en parte lo que ocurrió después: cuatro años entre la sentencia y la aprobación del Plan, con nuevas crecidas de por medio.

Este libro emplea ``ordena'' porque es el verbo que la propia ficha usa al describir el decisorio, y deja anotada la ambigüedad. La resuelve el texto íntegro.

\section{La aprobación del Plan y el apercibimiento penal}

\begin{ficha}{Audiencia del 12 de febrero de 2026}
Según el parte de prensa del Poder Judicial de Salta:
\begin{itemize}[leftmargin=*,nosep]
\item En audiencia encabezada por la jueza Mosmann, con intervención de las partes presentes y del \textbf{fiscal civil Agustín Vidal}, se resolvió \textbf{aprobar el Plan de Manejo Integral} presentado en autos, sin perjuicio del cumplimiento de las obligaciones legales y de las medidas ya establecidas en la normativa vigente.
\item Se dispuso \textbf{intimar al Municipio de La Caldera} a dar cumplimiento estricto, temprano y oportuno al \textbf{Protocolo de Crecidas e Inundaciones} establecido por ordenanza, a través del \textbf{Comité Operativo de Emergencia}, debiendo comunicar cada activación de modo inmediato al tribunal, \textbf{bajo apercibimiento de remitir los antecedentes a la justicia penal}.
\item Se establecieron plazos para que la Provincia presente la planificación de las obras y medidas prioritarias conforme el orden del Plan, y luego la del resto de las obras.
\item Se dispuso poner el texto del Plan, junto con el de la sentencia, a disposición de la comunidad en el domicilio constituido por el colectivo actor: Avenida General Güemes s/n, medidor 994, La Caldera.
\item Se destacó la labor del COPAIPA y de las profesionales que actuaron como veedoras, con nota de reconocimiento.
\end{itemize}
\end{ficha}

El apercibimiento penal es el dato nuevo y es el más severo de todo el expediente. Un tribunal civil que advierte a un municipio que remitirá los antecedentes a la justicia penal si no activa su propio protocolo de emergencia ante cada crecida está describiendo, sin decirlo, el nivel de confianza que le merece el cumplimiento voluntario después de siete años de litigio.

Conviene además retener la asimetría del decisorio. Al municipio se lo intima con apercibimiento penal a cumplir un protocolo. A la Provincia se le establecen plazos para presentar la planificación de obras. La diferencia de intensidad entre ambas órdenes reproduce, dentro del propio expediente, la distribución de responsabilidad que este libro describe: al que tiene menos capacidad se le exige más y con más consecuencia.

Y hay una tercera pieza, silenciosa. Que el texto del Plan se ponga a disposición de la comunidad en el domicilio constituido por el colectivo actor --- y no en la sede municipal, ni en el Boletín Oficial, ni en un sitio web --- indica que el tribunal consideró necesario asegurar el acceso por fuera de los canales institucionales. Es la misma cuestión que el capítulo \ref{cap:metodo} documentó para los anexos normativos, ahora resuelta por orden judicial.

\begin{cautela}
Una precisión de método antes de seguir, porque el verbo de la resolución importa. La jueza aprueba el Plan \textbf{presentado en autos}: es decir, aprueba un documento cuya elaboración estuvo a cargo de las propias demandadas. Que un plan haya sido aprobado no significa que resuelva la brecha entre el instrumento y su aplicación que este capítulo identifica. Verificarlo requiere leer el Plan contra el articulado del Código de Ordenamiento Territorial, que es la comprobación que este mismo capítulo enuncia como la más productiva.

Y hay un criterio de la condena de 2021 que sigue vigente y que conviene retener, porque es contrastable. La sentencia exigía una estrategia de estabilización basada en \textbf{infraestructura verde} --- reforestación, relocalización de elementos obstructores de la dinámica ecológica --- \textbf{por sobre la infraestructura gris} --- obras civiles, maquinaria pesada, edificación hormigonada ---, evaluando un balance progresivo de recuperación de la cuenca. Ese criterio se contrapone al plan de represas de retardo de crecida que el municipio propuso al jefe de Gabinete el 21 de enero de 2026 y que, según la propia comunicación municipal, fue aceptado. Contrastar ambos documentos es una línea abierta y estrictamente documental.
\end{cautela}

\section{El expediente digital: lo que el sistema judicial revela}

La consulta al sistema de expediente digital del Poder Judicial de Salta agrega tres elementos que ninguna otra fuente había aportado.

\begin{ficha}{Sistema de Expediente Digital --- CUIJ 17-00023860-0}
\begin{itemize}[leftmargin=*,nosep]
\item \textbf{Carátula principal:} \textit{Lazarte Vigabriel, Eduardo José Ignacio contra Municipalidad de La Caldera; Fiscalía de Estado por otros. \textbf{Tercero/s: Jardín Celestial S.R.L.}}
\item \textbf{Tribunal:} Juzgado de Minas, Secretaría Nº 2, Distrito Centro. Tribunal de segunda instancia: Corte de Justicia.
\item \textbf{Incidente:} INC 23860/1, ``piezas pertenecientes'', iniciado el 16 de marzo de 2022.
\item \textbf{Ejecución de sentencia:} Expte.\ Nº 780280.
\item \textbf{Tercero:} Jardín Celestial S.R.L., con apoderado José Luis Criado.
\item \textbf{Providencia del 22 de septiembre de 2022:} el tribunal tiene presente lo manifestado por el actor respecto de la \textbf{extemporaneidad e inadecuación del informe presentado por Fiscalía de Estado} y respecto de la \textbf{omisión por parte de la Municipalidad de Vaqueros}, y remite a la providencia dictada en el expediente de ejecución.
\item \textbf{Remisión a la Fiscalía Civil, Comercial y Laboral Nº 2:} 26 de marzo de 2024.
\item \textbf{Último movimiento del incidente a la fecha de esa consulta:} 16 de abril de 2026. La ejecución, en el Expte.\ 780280, registra actuaciones hasta septiembre (véase más abajo).
\end{itemize}
\end{ficha}

\begin{cautela}
El primero de los tres elementos es el más importante para este libro, porque une dos capítulos que hasta aquí corrían en paralelo.

\textbf{Jardín Celestial S.R.L.\ es tercero en el expediente.} Es la misma sociedad que, según la Resolución Delegada 369 D/2022, ofreció junto a Héctor D'Andrea\index{D'Andrea, Héctor} la donación de espacio institucional sobre la matrícula 6.558, correspondiente al proyecto urbanístico El Durazno, cuyo plano de mensura y loteo Nº 1.164 fue aprobado por Resolución Delegada 646/21.

Es decir: el desarrollador del loteo que se inundó está incorporado al proceso judicial que litiga por esa inundación. La ficha académica del caso consignaba que la sentencia ``no implica obligaciones para terceros, sin perjuicio de las cuestiones que puedan surgir entre estos durante la etapa ejecutiva''. El sistema muestra que un tercero fue efectivamente incorporado.

Nada permite afirmar en qué carácter ni a instancia de quién --- la citación de tercero puede ser voluntaria o forzosa, y el registro consultado no lo consigna. Lo que en cambio queda establecido, con fuente judicial, es que la sociedad titular de El Durazno es parte del expediente. Determinar cuándo se la citó, por qué motivo y con qué resultado es materia del expediente completo, y es probablemente la pregunta más productiva que quedó abierta en todo este libro.
\end{cautela}

El segundo elemento corrige el alcance territorial del caso. Este libro trató el amparo como un asunto exclusivamente caldereño. La providencia del 22 de septiembre de 2022 registra que \textbf{la Municipalidad de Vaqueros fue requerida y omitió informar}.

Vaqueros no figura como demandada en la carátula, pero fue alcanzada por un requerimiento del tribunal y no respondió. Eso tiene sentido hidrológico --- la microcuenca del río La Caldera desemboca en el sistema que atraviesa Vaqueros, y el río Wierna marca el límite entre ambos municipios, conforme la Ley 4987 de 1974 (capítulo \ref{cap:cot}) --- y tiene una consecuencia para el capítulo \ref{cap:vaqueros}: el municipio cuyo Código, en su texto de 2026, define un área de protección hídrica como urbanización parque diferida ya había sido requerido, cuatro años antes, en un litigio sobre el manejo de esa cuenca.

En la misma providencia el actor cuestiona la \textbf{extemporaneidad e inadecuación} del informe de la Fiscalía de Estado. El triángulo de imputaciones que este capítulo describe más abajo no se resolvió con la sentencia: continuó durante la ejecución, ahora bajo la forma de informes tardíos e informes que no se presentan.

\begin{cautela}
El tercer elemento es de escala temporal. La causa se inicia con hechos de diciembre de 2018, obtiene sentencia en diciembre de 2021, abre un incidente en marzo de 2022, registra omisiones de informe en septiembre de 2022, pasa a la Fiscalía Civil en marzo de 2024, aprueba el Plan en febrero de 2026 y sigue con movimientos en abril de ese año.

Siete años y cuatro meses desde los hechos hasta abril de 2026 ---y la causa sigue---, con nuevas crecidas de por medio, para un amparo --- que es, por definición legal, la vía expedita y rápida del artículo 43 de la Constitución Nacional.

La segunda instancia consignada es la Corte de Justicia de la Provincia, y la causa llegó allí: los pedidos municipales de fondos de 2025 y 2026 consignan costas impuestas por la Corte en forma solidaria a la Provincia y al municipio. Lo que la Corte resolvió sobre el fondo sigue sin leerse.
\end{cautela}

Una precisión sobre la carátula. La ficha académica del Observatorio consigna como demandados a la Municipalidad de La Caldera y al Ministerio de la Producción, Trabajo y Desarrollo Sustentable; el sistema judicial consigna a la Municipalidad de La Caldera y a la \textbf{Fiscalía de Estado}. No hay contradicción: la Fiscalía de Estado es el órgano que representa a la Provincia en juicio. Pero la diferencia importa para quien busque el expediente, y por eso se consigna el CUIJ, que es el identificador unívoco: \textbf{17-00023860-0}.

\textbf{Y ahora la carátula puede leerse en el propio expediente.} Las dos providencias que este apartado citaba por el sistema digital --- la del 16 de agosto y la del 22 de septiembre de 2022 --- fueron obtenidas en copia. La primera transcribe la carátula del principal: «\textit{Eduardo Jose Ignacio Lazarte Vigabriel vs.\ Municipalidad de La Caldera y Ministerio de Producción, Trabajo y Desarrollo Sustentable de la Provincia de Salta --- Amparo Colectivo. Expte.\ Nº MIN 23860/19}». \textbf{El demandado es el Ministerio}, como consignaba la ficha académica; la Fiscalía de Estado aparece en el sistema porque es quien lo representa. La duda que este apartado dejaba abierta queda cerrada, y a favor de la fuente académica.

Esa misma providencia agrega el dato de trámite que explica por qué el material está donde está: como el principal se encuentra radicado en la Corte de Justicia, el juzgado dispuso que \textbf{las presentaciones de la etapa de cumplimiento se hagan en las piezas individualizadas como INC-23860}. Todo lo que las demandadas acompañaron en cumplimiento de la condena --- los mapas, la nómina de canteras, los planos del Código --- fue a parar a un incidente, no al expediente principal. \textbf{Es la razón material por la que esta cartografía no está en ningún lado donde alguien la busque.}

\section{Las defensas}

La parte más valiosa del expediente no es la sentencia: son las contestaciones.

Citada por la justicia, la Secretaría de Recursos Hídricos provincial pidió el rechazo del amparo cuestionando la legitimación del actor, sostuvo que el organismo no tenía capacidad para hacer obras y apuntó contra los responsables de los loteos como encargados de construir las defensas. El municipio adhirió a esa postura. La Fiscalía de Estado, en cambio, indicó que en la medida en que los loteos son autorizados por el municipio, es el municipio quien debe hacer las obras y garantizar las condiciones.

Léase la secuencia completa. La provincia dice que la responsabilidad es de los loteos. El municipio adhiere a que la responsabilidad es de los loteos. La Fiscalía de Estado --- órgano de la propia provincia --- dice que la responsabilidad es del municipio, porque el municipio autoriza.

El resultado agregado de las tres posiciones es que nadie es responsable y que el daño ya ocurrió. Y sin embargo cada una, tomada por separado, es una posición procesal defendible.

Este triángulo es la mejor pieza documental del libro, y es mejor que cualquier cosa que una investigación externa hubiese podido reconstruir, por la misma razón que hacía valiosa, en \textit{El dispositivo salteño}, la frase de Saravia Frías en el Club 20 de Febrero: el dispositivo se describe con exactitud cuando se defiende, porque al defenderse tiene que explicar cómo funciona.

\subsection*{Lo que el propio Código ya había dicho}

El capítulo \ref{cap:cot} documentó que el artículo 42 del Código --- el mismo instrumento al que la sentencia manda ajustar el Plan --- clasifica expresamente a El Nogalar, El Durazno y Potreros de La Caldera como núcleos urbanos espontáneos o no planificados, con alumbrado y recolección en vías principales como máximo.

Ese artículo cambia el sentido procesal del expediente. Cuando el municipio adhirió a la postura de que la responsabilidad correspondía a los titulares de los loteos, ya había declarado por ordenanza que esos loteos eran no planificados y que él mismo les proveería servicios mínimos. La defensa municipal y su propia norma de ordenamiento dicen, sobre el mismo objeto, cosas compatibles entre sí pero devastadoras en conjunto.

\section{Una clasificación que reordena el análisis}

Una nota de \textit{El Tribuno} de enero de 2026 aporta un dato central: los loteos El Nogalar 1 a 4 son \textbf{loteos municipales}. Los demás son emprendimientos privados, y Potreros de La Caldera es un barrio cerrado ubicado muy cerca del río. La misma nota consigna que, según fuentes municipales, esos loteos son irregulares; que según el intendente las propiedades de esos barrios no pagan impuestos municipales; y que, para el jefe comunal, los emprendimientos privados le piden asistencia al Estado.

Si El Nogalar es de origen municipal, el municipio no es solamente el organismo que autorizó loteos privados defectuosos: es también el loteador. La pregunta deja de ser por qué no controló al mercado y pasa a ser en qué condiciones técnicas produjo suelo urbano por cuenta propia sobre una microcuenca cuya fragilidad estaba advertida, al menos, desde el plan urbano de 2015. En enero de 2026 el intendente declaró a la prensa que El Nogalar 4 se encuentra dentro del lecho del río, carece de catastro y fue desarrollado durante gestiones municipales anteriores, y que muchos barrios se asentaron en zonas inundables autorizadas, según él, por informes técnicos de Recursos Hídricos. Son declaraciones y no documentos, pero la falta de catastro es verificable.

La respuesta más probable no es venal. Un municipio pequeño que necesita resolver demanda habitacional y no dispone de capacidad técnica hidrológica lotea donde puede y con lo que tiene. Pero esa explicación, si es la correcta, tiene una consecuencia política que el libro nombra en el capítulo \ref{cap:hacienda}: la insuficiencia técnica municipal no es un accidente administrativo, es el modo en que la provincia distribuye responsabilidad hacia abajo sin distribuir capacidad.

\section{Quién escribió el Plan}
El Plan de Manejo Integral que la jueza aprobó en febrero de 2026 no lo redactaron los organismos técnicos del Estado condenado. Lo redactó una consultora privada, contratada por la Provincia dos años después de la condena, sin licitación.

\begin{ficha}{Contratación Abreviada --- Boletín Oficial Nº 21.640, 26 de enero de 2024}
\textbf{Expediente Nº 136-221704/23.} Ministerio de Producción y Desarrollo Sustentable.

\textbf{Objeto:} contratación del servicio de \textbf{consultoría externa} para la elaboración y presentación del plan de manejo integral de la microcuenca del río La Caldera, ``en cumplimiento a la demanda judicial dictada'' en los autos \textit{Eduardo José Ignacio Lazarte Vigabriel contra la Municipalidad de La Caldera contra Fiscalía de Estado, Municipalidad de La Caldera \textbf{por otros tercero/s Jardín Celestial S.R.L.}}, \textbf{Expediente Judicial Nº 238670/20}.

\textbf{Destino:} Secretaría de Recursos Hídricos. \textbf{Proveedor: YSATI INGENIERÍA S.R.L.}

\textbf{Encuadre:} Ley Nº 8.072, contratación abreviada, \textbf{artículo 15 inciso e)}. Salta, 23 de enero de 2024. Firma: Dorigato Manero, Coordinador General. Orden de publicación 100111211.
\end{ficha}

\begin{cautela}
Cuatro puntos que este aviso establece y que ninguna otra fuente de este capítulo establecía.

\textbf{Primera: la carátula oficial nombra a Jardín Celestial S.R.L. como tercero.} Hasta aquí la vinculación entre la sociedad y la causa se apoyaba en el expediente y en la ficha del Observatorio. Ahora consta en el Boletín Oficial, en un acto administrativo del Poder Ejecutivo provincial, sin intermediación periodística. Y aparece un número de expediente judicial que este libro no tenía: \textbf{238670/20}, distinto del 780280 que registra la audiencia de 2026. Muy probablemente es el EMI 23860/2020 con un dígito de más, error de tipeo del aviso; conviene cotejarlo antes de tratarlo como un expediente distinto.

\textbf{Segunda: el plan se tercerizó.} La condena de diciembre de 2021 obligó a la Provincia y al Municipio a presentar un plan de manejo. La Provincia cumplió comprando el plan. El destino es la Secretaría de Recursos Hídricos --- el mismo organismo que había fijado la línea de ribera de los arroyos El Durazno y Guaranguay en diciembre de 2014, y que en el amparo pidió el rechazo de la demanda sosteniendo que no tenía capacidad para hacer obras. El aviso, leído junto a esa defensa, es coherente: el organismo declaró no tener capacidad y después contrató la capacidad afuera.

\textbf{Tercera: el encuadre invocado no es la urgencia.} El artículo 15 de la Ley 8072 tiene un inciso para razones de urgencia --- el i) ---, que es el que la Provincia usa habitualmente en obra hídrica del Valle de Lerma. Aquí se invocó el \textbf{inciso e)}, cuya reglamentación exige fundar la capacidad científica, técnica, tecnológica, profesional o artística de la persona o empresa a la que se encomienda la prestación. Es decir: se contrató sin competencia por especialidad, no por apuro. Dos años después de una condena judicial, el apuro ya no era invocable.

\textbf{Cuarta: los tiempos.} Condena en diciembre de 2021. Contratación de la consultoría en enero de 2024, \textbf{veinticinco meses después}. Aprobación judicial del plan en febrero de 2026, \textbf{cincuenta meses después de la condena}. Entre el segundo aluvión y la aprobación del plan que debía prevenirlo transcurrieron siete años y un mes.
\end{cautela}

\begin{cautela}
Esto obliga a matizar, no a retirar, la advertencia de método formulada más arriba. Se dijo allí que la jueza aprobó un documento elaborado por las demandadas. Es más preciso decir que aprobó un documento \textbf{elaborado por una consultora contratada por una de las demandadas}, mediante un procedimiento sin competencia, y sometido a su vez a la veeduría del COPAIPA en una etapa del proceso.

Nada de esto pone en duda la calidad técnica del Plan, que este libro leyó pero no está en condiciones de evaluar como obra de ingeniería. Lo afirmable es que quedó fuera de la contratación, por definición, la parte que la condena de 2021 no podía comprar: la capacidad estatal permanente de gestionar la cuenca. Un plan se encarga; una capacidad, no.
\end{cautela}

\section{Los mapas de la cuenca existen desde 2022}

Cualquier plan para la cuenca necesita, antes que nada, su cartografía. Ya
existía, y se produjo por mandato de esta sentencia.

\begin{ficha}{«Informe de mapas de cuenca de río La Caldera» --- Secretaría de Recursos Hídricos, mayo de 2022}
Ocho hojas con \textbf{tres mapas temáticos}, a escalas \textbf{1:125.000} y \textbf{1:40.000},
en sistema de referencia \textbf{POSGAR 2007 --- Faja 3}.

Firman el Ing.\ Guillermo A.\ Fuchs, de la Dirección de Fiscalización y Control; el Ing.\ Civil
Ricardo Cristóbal Ramos, Jefe de Programa Proyectos Hídricos; y Mauricio Romero Leal.

El informe delimita la cuenca del río La Caldera dentro de la alta cuenca del Mojotoro ---
que a su vez integra la del Bermejo --- e individualiza sus \textbf{subcuencas}, con la
divisoria de aguas descripta por cotas y topónimos: cerro 4372, abra Los Yacones, cerro
Vaqueros (1.843 m), serranías de Lesser y de Mojotoro.
\end{ficha}

\textbf{De modo que la cuenca está delimitada oficialmente, a escala útil y en un sistema de
coordenadas estándar, desde mayo de 2022.} Escala 1:40.000 permite ubicar parcelas; POSGAR
2007 es el marco de referencia oficial argentino, de modo que esa cartografía \textbf{es
superponible con cualquier plano catastral}.

El informe individualiza además las subcuencas con su superficie y su perímetro, y ese cuadro
es la primera medición oficial del territorio que este libro encuentra desagregada por curso.

\needspace{9\baselineskip}
{\sffamily\bfseries\footnotesize Superficies de la cuenca y sus subcuencas, según el informe de mayo de 2022\par}
\vspace{0.5em}
{\small
\setlength{\LTpre}{0.6em}\setlength{\LTpost}{1.2em}
\begin{longtable}{@{}p{0.9cm}p{6.4cm}r r@{}}
\toprule
\textbf{fid} & \textbf{Nombre} & \textbf{Sup.\ (ha)} & \textbf{Perím.\ (m)} \\
\midrule
\endhead
1 & Analuya & 2.072,95 & 30.217 \\
2 & Cuchimayo & 2.323,50 & 37.084 \\
4 & \textbf{El Durazno} & \textbf{1.391,96} & 27.644 \\
5 & Finca de Linares --- Guaranguay & 919,09 & 20.828 \\
7 & La Angostura & 1.626,09 & 28.504 \\
8 & La Caldera & 6.533,71 & 92.270 \\
9 & La Calderilla & 2.018,78 & 27.975 \\
14 & San Alejo & 3.948,67 & 53.417 \\
15 & Santa Rufina & 2.935,84 & 45.975 \\
\bottomrule
\end{longtable}
}

\begin{cautela}
El cuadro se transcribe tal como consta. \textbf{La numeración tiene huecos}: faltan los
identificadores 3, 6 y 10 a 13, de modo que el informe trabajó sobre una capa de al menos
quince polígonos y elevó nueve. No puede saberse desde afuera si los ausentes quedaron fuera
del área de interés o si simplemente no se listaron, y la pregunta corresponde hacerla junto
con el pedido de las capas vectoriales.

\textbf{Y hay una segunda partición de la misma cuenca, en el mismo expediente, que no coincide
con ésta.} La presentación del Plan Maestro de Drenaje que este capítulo analiza más abajo
---lámina \ref{lam:pmdp}--- enumera \textbf{siete} subcuencas: Volcán--Yacones, Santa Rufina,
San Alejo, Las Nieves, Lesser, Castellanos y Angosto. El cuadro de mayo de 2022 enumera
\textbf{nueve}, y de los dieciséis nombres sólo dos se repiten ---San Alejo y Santa Rufina---,
más «La Angostura», que podría ser el «Angosto» de la otra. Analuya, Cuchimayo, El Durazno,
Finca de Linares--Guaranguay, La Caldera y La Calderilla no figuran en la segunda; Volcán--Yacones,
Las Nieves, Lesser y Castellanos no figuran en la primera. \textbf{Dos documentos del mismo
organismo, sobre la misma cuenca y en la misma pieza judicial, la dividen de dos maneras que no
son traducibles una a la otra.} Pueden corresponder a dos escalas o a dos objetos distintos ---una
al área del Plan de Manejo, otra al área del estudio de drenaje---, y este libro no lo establece:
lo dirían las capas vectoriales, que sigue pidiendo. Lo que registra es que ninguno de los dos
documentos advierte de la existencia del otro.
\end{cautela}

Dos consecuencias, y una tercera que este apartado debe corregir sobre sí mismo.

\textbf{La primera corrige una carencia que este libro declaraba en versiones anteriores.} La delimitación de la cuenca no falta: se produjo hace cuatro años, por
orden judicial, y está firmada por tres profesionales del organismo hídrico.

\textbf{La segunda es dónde están esos mapas.} No en un boletín oficial ni en un portal
público, sino en una pieza de un expediente judicial. \textbf{Existen, son oficiales, y para
llegar a ellos hay que saber que existen y en qué pieza del expediente están.} Es la misma
mecánica que el capítulo \ref{cap:opacidad} describe: el acto se publica, el contenido que permitiría
verificarlo queda en otro lado.

\textbf{Y la tercera: este relevamiento ya los tiene.} Una versión anterior de este apartado
decía que no. Las ocho hojas se incorporaron, y cuatro de ellas se reproducen en este libro: el
mapa de la cuenca y sus microcuencas (lámina \ref{lam:mapa1}), el de las líneas de ribera
registradas, que acompaña al capítulo \ref{cap:ribera} (lámina \ref{lam:mapa2}), y el de los polígonos
mineros, que acompaña al capítulo \ref{cap:resistencias} (lámina \ref{lam:mapa3}). La hoja 1 del informe
--- con la cuenca del Mojotoro y las tres firmas --- va en la lámina \ref{lam:mojotoro}.

\begin{cautela}
Conviene decir con qué se hicieron. La última hoja del informe declara los materiales:
\textbf{QGIS 3.16.7}, modelo de elevación \textbf{SRTM}, imágenes de base de Google Earth y
capas del IGN, del geoportal \textbf{IDESA} y de la propia Secretaría. Y declara las fuentes,
entre las cuales hay dos estudios de la Secretaría de Recursos Hídricos --- de 2008 y de 2009,
sobre cuencas y sobre el Mojotoro --- consignados expresamente como \textbf{«inédito»}.

\textbf{La cartografía oficial de esta cuenca se apoya, en parte, en documentos que el propio
organismo declara no publicados.} No es una irregularidad y no se la presenta como tal: es la
tercera capa de la misma estructura que este libro viene describiendo.

Y un detalle de soporte, que vale para las cuatro láminas de origen hídrico. La hoja 2 del
informe lleva por epígrafe «Mapa de Cuenca de río La Caldera y \textbf{en color} cuencas de
áreas de interés». \textbf{El ejemplar incorporado al expediente es en escala de grises.} El
color, si existió, no está en la pieza que el juzgado recibió, y con él se pierde la distinción
que el epígrafe anuncia.
\end{cautela}

\lamina{lam-mapa1-cuenca.jpg}{lam:mapa1}%
{Cuenca hidrográfica del río La Caldera y microcuencas de sus cauces aportantes}%
{\textbf{Mapa Nº 1.} Cuenca hidrográfica del río La Caldera y las microcuencas de sus cauces
naturales aportantes. Escala 1:125.000, sistema POSGAR 2007 --- Faja 3. Es la cartografía con
la que se estableció la superficie del Plan de Manejo Integral: la cuenca media y los arroyos
con desembocadura en el afluente principal.}%
{«Informe de mapas de cuenca de río La Caldera», Secretaría de Recursos Hídricos de Salta, mayo
de 2022, presentado en la ejecución de sentencia del amparo colectivo, expediente
INC-23860/1. Reproducción facsimilar; se recortó el margen en blanco del escaneo y no se
alteraron los tonos.}

\lamina{lam-mojotoro-firmas.jpg}{lam:mojotoro}%
{Cuenca del río Mojotoro y cuenca del río La Caldera, con las firmas del informe}%
{\textbf{Hoja de resumen del informe}, con la delimitación de la alta cuenca del Mojotoro
dentro de la cual se individualiza la del río La Caldera, y el pie de firmas: el Ing.\
Guillermo A.\ Fuchs, de la Dirección de Fiscalización y Control; el Ing.\ Civil Ricardo Cristóbal Ramos, Jefe de Programa Proyectos Hídricos; y Mauricio Romero Leal, cuyo sello está cortado
y velado por la rúbrica: se lee el nombre y «\ldots\ de Recursos Hídricos», no el cargo
completo.}%
{«Informe de mapas de cuenca de río La Caldera», Secretaría de Recursos Hídricos de Salta, mayo
de 2022, hoja 1, expediente INC-23860/1. Reproducción facsimilar; se recortó el margen en
blanco del escaneo y no se alteraron los tonos.}

\section{El Plan de Manejo}

Este libro reconstruyó durante meses un documento que no había leído. Ya no es así. El Plan
Integral de Manejo de la Microcuenca del Río La Caldera son doscientas treinta y cuatro hojas
y doce planos, y lo que dice obliga a revisar varias cosas.

\subsection*{Quién lo escribió}

Lo presentó el ingeniero \textbf{Francisco Javier Ramos Vernieri por YSATI Ingeniería
S.R.L.}, en el expediente 136-221704/23, bajo la carátula «Plan Director de la Municipalidad
de La Caldera». La Dirección General de Infraestructura Hídrica no formuló objeciones el 14
de noviembre de 2024, y Fiscalía de Estado lo acompañó al juzgado tres semanas después.

\textbf{Es la misma firma a la que este libro dirige uno de sus pedidos de derecho a
réplica}, preparados aparte y no incluidos en este volumen. La carátula del expediente ---«Plan Director de la Municipalidad de La Caldera»--- presenta el trabajo también como instrumento de planificación municipal; qué alcance tiene esa denominación, y si la firma intervino en otros instrumentos del municipio, es una de las preguntas de ese pedido.

\subsection*{El diagnóstico social, y lo que revela sobre este libro}

El capítulo más extenso del informe --- setenta y cinco páginas --- es el diagnóstico social,
y se apoya en \textbf{encuestas estructuradas y semiestructuradas} realizadas entre el 14 de
febrero y el 18 de julio de 2024, en papel y en formato digital.

La metodología explica por qué se eligió así, y la explicación merece citarse:

\begin{cautela}
Se usó formato mixto para abarcar más barrios y más horarios «ya que la mayoría son casas de
campo o trabajan en la ciudad regresando a sus hogares al final del día» --- lo que confirma,
desde adentro, lo que el capítulo \ref{cap:poblacion} mide con los censos ---, y \textbf{para «ofrecer el
anonimato para las personas que expresaron su temor a “represalias” o “consecuencias
negativas” por brindar su opinión»}.

Los encuestadores agregan una razón: «la ejecución de la sentencia llevó varios años en
proceso, lo que generó grandes tensiones entre los distintos grupos sociales».

\textbf{Este libro declara, entre sus límites, que no contiene entrevistas propias} ---apenas dos conversaciones informales con vecinos, que el capítulo \ref{cap:poblacion} registra como testimonio---. El
Plan contiene decenas, y sus formularios en papel quedaron «en resguardo de la consultora y a
disposición del juzgado». Es el material que a este trabajo le falta, y existe.
\end{cautela}

El planteo del capítulo social coincide con lo que este libro documenta por vía documental, y
conviene verlo dicho por un técnico contratado por el propio Estado:

\begin{ficha}{Plan de Manejo, capítulo 10, página 100}
«En los últimos 20 años, la localidad de la Caldera ha tenido una expansión urbana hacia el
sur, la cual \textbf{carece de infraestructura vial, pluvial y de defensa fluvial}.»

«La falta de planificación urbana, y la gran demanda de terrenos en la zona generan
\textbf{asentamientos privados y edificaciones espontáneas en lugares poco convenientes}.»

«La comunidad está formada por los pobladores locales y por los nuevos habitantes con
distintas actividades económicas y culturales, que generan \textbf{tensiones sociales
agravadas por la falta de espacios de diálogo y la influencia de la política partidaria}.»
\end{ficha}

\subsection*{Quiénes contestaron, y quiénes no}

El Plan publica el perfil de la muestra, y es el retrato de población más detallado que este
libro consiguió.

\textbf{Cuarenta y ocho encuestas efectivas, sobre ciento setenta y cuatro personas} relevadas. (El informe del taller de 2023 también registra cuarenta y ocho encuestas completadas; que la coincidencia no sea una confusión de series es algo que conviene cotejar.) El setenta y siete por ciento de los encuestados es jefe de hogar, y de ésos
\textbf{el cincuenta y ocho por ciento son mujeres}. Las edades van de dieciocho a ochenta y
dos años. El cincuenta y dos por ciento es trabajador independiente, emprendedor o artesano;
el treinta y tres por ciento, ama de casa.

Sobre la vivienda: el noventa y dos por ciento se usa únicamente como vivienda --- no como
casa de fin de semana ---, el ochenta y ocho por ciento tiene servicio eléctrico «con
conexiones precarias a medida que más se alejan del ejido urbano», y sobre el agua:
\textbf{«la mayoría de las casas debe recolectarla en tanques y/o cisternas, teniendo en
cuenta que casi el 50\,\% de los encuestados pertenecen a grupos vulnerables (menores de edad
y adultos mayores)»}.

\textbf{Y once personas se negaron a contestar.} El Plan dice por qué, y hay que citarlo
entero:

\begin{cautela}
«Fueron 11 las personas que decidieron no contestar las encuestas, ya que, al momento de
comentarles el objetivo del relevamiento, adujeron \textbf{no tener “esos problemas” y “no
querer tener problemas con nadie”}; pertenecían a los barrios: Santa Mónica, Zavaleta y
Balcón de Getsemaní; que se encuentran dentro del área de influencia indirecta.»

Las dos razones son distintas y conviene separarlas. \textbf{«No tener esos problemas»} es la
descripción exacta de lo que el capítulo \ref{cap:poblacion} mide: quien vive fuera del área de riesgo, o no
vive allí todo el año, no comparte el problema del río. \textbf{«No querer tener problemas
con nadie»} es otra cosa, y es la misma que llevó a la consultora a ofrecer anonimato.

Uno de los tres barrios se llama \textbf{Balcón de Getsemaní}. «Getsemaní» es el nombre de la
finca que figura como domicilio del gerente titular ---desde 2018; socio desde 2010--- de la sociedad titular del loteo El
Durazno, según documenta el capítulo \ref{cap:loteo}.
\end{cautela}

\subsection*{Las conclusiones del diagnóstico}

La página 184 enumera lo que encontró. Son cinco puntos y cada uno confirma un capítulo de
este libro.

\textbf{«Poca participación ciudadana \ldots\ Las acciones son direccionadas desde el municipio
sin un trabajo de concientización, escucha y participación de los vecinos.»}

\textbf{«Deficientes trabajos de encauzamiento de ríos y arroyos \ldots\ En general estos
trabajos son realizados por el municipio y por privados \textit{sin autorización de la
Autoridad de aplicación}.»} El capítulo \ref{cap:defensas} documenta ciento dieciséis años de obras sobre ese cauce; el Plan agrega que algunas se hacen \textbf{sin autorización}.

\textbf{«Incorporación de urbanizaciones en sectores de riesgo fluvial y pluvial \ldots\ El
municipio no cuenta con recursos técnicos ni específicos para llevar adelante estas tareas.»}
Es, palabra por palabra, la tesis de incapacidad técnica que la Advertencia preliminar de
este libro adopta como regla de lectura.

\textbf{«Los controles \ldots\ son deficientes y suelen observarse en general cuando ya los
hechos se encuentran consumados y carecen de seguimiento e implementación de acciones de
reparación del daño generado.»}

\textbf{Y la falta de infraestructura de drenaje pluvial, que «debe ser solicitado con suma
rigurosidad a los desarrolladores inmobiliarios».}

\subsection*{Lo que cuesta}

El Plan cifra sus medidas estructurales. A precios de septiembre de 2024:

\needspace{8\baselineskip}
{\small
\begin{longtable}{@{}p{6.2cm}p{1.6cm}p{4.2cm}@{}}
\toprule
\textbf{Medida} & \textbf{Tipo} & \textbf{Inversión} \\
\midrule
\endhead
Defensa río La Caldera & fluvial & \$13.696.779.462 \\
Camino de ribera & fluvial & \$2.390.083.014 \\
Defensa río Guaranguay & fluvial & \$2.330.019.749 \\
Camino secundario & fluvial & \$1.865.431.693 \\
Defensa río Nogalar & fluvial & \$1.426.097.368 \\
Defensa río Durazno & fluvial & \$892.825.605 \\
Bosques de ribera & fluvial & \$226.114.410 \\
Seis canales pluviales urbanos & urbana & \$972.834.733 \\
\addlinespace
\multicolumn{3}{@{}p{12cm}@{}}{\footnotesize\textit{Los rubros se transcriben con la nomenclatura del Plan. En el resto de este libro, El Durazno y Guaranguay son \textbf{arroyos} ---así los nombran las Resoluciones 328/12 y 383/14--- y «Nogalar» no es un curso de agua sino un loteo municipal: el único arroyo de nombre parecido que registra la serie es el \textbf{Los Nogales Norte o El Cajón} de la Resolución 101/14, que atraviesa la Finca Vaqueros, y su relación con el sector no está establecida (capítulo \ref{cap:ribera}). Si el «río Nogalar» del Plan es ese arroyo o es otro, el Plan no lo dice.}} \\
\midrule
\textbf{Total} & & \textbf{\$23.800.186.036} \\
\bottomrule
\end{longtable}
}

\begin{cautela}
Las cifras se cotejaron contra la imagen de la tabla 7 del Plan (página 221; lámina \ref{lam:pmtabla7}) y de la tabla 8 (páginas 222 y 223), que coinciden entre sí. \textbf{La suma de los renglones cierra al centavo con el total impreso}, \$23.800.186.036,87 ---con los centavos del original; sin ellos, como se transcriben en la tabla, los renglones suman \$23.800.186.034---, y lo mismo ocurre con la operación anual, \$11.900.093,02. Una versión anterior de este cuadro, tomada de un reconocimiento óptico, daba un total de \$23.600 millones y un desajuste entre renglones y total; los dos errores eran de la lectura, no del Plan.
\end{cautela}

\lamina{lam-pm-tabla7.jpg}{lam:pmtabla7}%
{Inversión y mantenimiento anual por medida en el Plan de Manejo}%
{\textbf{Tabla 7 del Plan de Manejo Integral de la Microcuenca del Río La Caldera: «Costo anual de mantenimiento».} CAPEX y OPEX de las trece medidas estructurales, a precios de septiembre de 2024. Los renglones suman exactamente los totales impresos, \$23.800.186.036,87 y \$11.900.093,02; en cada fila, el OPEX es el 0,05\,\% del CAPEX.}%
{Expediente 0090136-221704/2023-2, Secretaría de Recursos Hídricos de Salta; informe de YSATI Ingeniería S.R.L., emitido el 31/10/2024, página 221. Recorte del escaneo; no se alteró la imagen.}

Tres lecturas.

\textbf{La defensa del río La Caldera concentra el cincuenta y ocho por ciento de la
inversión.} Trece mil setecientos millones de pesos sobre veintitrés mil ochocientos. Es, en su escala, la obra que el decreto de enero de 1935 llamó «la solución definitiva» y que en noventa años nunca se ejecutó. \textbf{Este Plan es la primera vez que aparece proyectada y cifrada.}

\textbf{Y hay una línea que vale por sí sola: «bosques de ribera».} Doscientos veintiséis
millones, menos del uno por ciento del total. Es la única medida de \textbf{infraestructura
verde} del cuadro --- exactamente lo que la sentencia de 2021 ordenó privilegiar por sobre la
gris.

\textbf{Y la tercera es sobre el mantenimiento.} El Plan estima la operación y el mantenimiento anual de las trece obras en \$11,9 millones, y lo define como «la inversión municipal para conservar las obras de desagües urbanos y defensas», con un horizonte de veinte años. Pero la cifra no sale de un cálculo obra por obra: \textbf{en las trece medidas, el costo anual de mantenimiento es exactamente el 0,05\,\% de la inversión}. En los veinte años del horizonte, eso suma el 1\,\% de lo invertido, a cargo de un municipio de tercera categoría (capítulo \ref{cap:poblacion}).

\subsection*{La alerta temprana existe, y está proyectada}

El Plan propone cinco estaciones: \textbf{San Alejo, Santa Rufina y Azud Campo Alegre}
hidrométricas, \textbf{Alta Caldera y Municipio} meteorológicas, todas con coordenadas.

Y explica por qué funcionaría: el punto de medición está «a 5 Km aguas arriba de la localidad
de La Caldera, lo que implica que aproximadamente el evento demora \textbf{entre 45 min a 60
min en llegar a la trama urbana, siendo un tiempo suficiente para alertar a la población y a
los organismos de emergencias}».

\textbf{Cuarenta y cinco minutos.} Eso es lo que separa a ese pueblo de una crecida, y lo que
un sistema de alerta le devolvería. En la audiencia de febrero de 2026, el actor manifestó
que el Servicio Meteorológico Nacional le había informado que esos dispositivos «requieren de
inversión mínima». Dos de las cinco estaciones propuestas llevan los nombres de fincas que
este libro persigue desde 1909: San Alejo y Santa Rufina.

\section{Lo que se dijeron en la audiencia}

El acta de la audiencia del 12 de febrero de 2026 transcribe lo que cada parte sostuvo antes
de que la jueza resolviera. Es la única pieza de este libro donde el Estado explica, en sus
propios términos, por qué cuatro años después de una condena no había obras.

\textbf{El actor abre diciendo que en 2025 no hubo avances.} El Dr.\ Chávez manifiesta que
durante ese año «no se observó avances en las acciones aprobadas», que en la última época de
lluvias ocurrieron inundaciones \textbf{«que ya fueron previstos en el Plan presentado en
2024»} sin que la Provincia ni el municipio adoptaran medidas, y que quiere saber por qué no
se hizo lo que los propios condenados iban a hacer.

\textbf{La Provincia contesta con una distinción jurídica, y hay que citarla completa.} La
Dra.\ Herrando manifiesta que su parte estaba condenada a presentar un Plan de Manejo, que lo
presentó en tiempo y forma, y que \textbf{«la presentación del Plan no implicaba que la
Provincia iba a realizar obras»}: primero hay que determinar cuáles, prever la partida
presupuestaria y cumplir los procedimientos administrativos. Y agrega que «la sentencia no
fija que deberían haberse ejecutado de inmediato obras sino primero la presentación del
Plan».

\textbf{Y la Subsecretaría de Recursos Hídricos lo dice todavía más claro.} El Dr.\ Le Favi
expresa que «se considera necesario que \textbf{la aprobación del Plan sea el catalizador de
partidas presupuestarias}».

Es decir: \textbf{el Estado sostiene que no podía presupuestar obras hasta que un juez
aprobara el Plan}, y el Plan estuvo cuatrocientos treinta y tres días en el juzgado esperando esa
aprobación. Entre la presentación y la aprobación hubo dos temporadas de lluvias.

\begin{cautela}
La posición de la Provincia es jurídicamente defendible y este libro no la descalifica: una
condena a presentar un plan no es una condena a ejecutar obras, y el gasto público exige
partida previa. El Dr.\ Chávez sostiene lo contrario --- que la segunda etapa de la sentencia
obligaba a comenzar a ejecutar y a prever las partidas --- y \textbf{ese desacuerdo es sobre
el alcance de la sentencia de 2021, cuyo texto este libro no leyó}.

Lo que sí queda establecido, porque lo dicen las dos partes, es la \textbf{secuencia real}:
sentencia en 2021, Plan presentado en diciembre de 2024, aprobado en febrero de 2026, y la
planificación de obras recién comprometida para cuarenta y cinco días después de esa fecha.
\end{cautela}

\subsection*{Cinco cosas que el acta revela y que este libro no tenía}

\textbf{Existe un ``Plan de Mínimas'', distinto del Plan de Manejo, y es muy anterior al
juicio.} Lo menciona dos veces la Subsecretaría en la audiencia de 2026 como algo elaborado
«independientemente». Dos fuentes permiten fecharlo y describirlo.

En noviembre de \textbf{2010} ya se ejecutaban trabajos en el departamento bajo el ``Plan de
Mínima de Defensas y Encauzamientos de Ríos'': treinta y ocho municipios, diez mil horas de
máquina y cuarenta kilómetros de defensa sobre cuarenta y cinco ríos de la provincia. Ese año
se intervino, con doscientas cincuenta horas de máquina, en el río La Caldera y en los
arroyos \textbf{Guaranguay, Los Duraznos y Los Manzanos}, para proteger el camping y el barrio
El Jardín.

Y en \textbf{2016} el instrumento se formalizó por decreto, con anexo publicado.

\begin{ficha}{Decreto Nº 1.889/16 --- Acta Acuerdo, 23 de septiembre de 2016}
«Plan Bicentenario de la Independencia --- \textbf{Plan de Mínimas y Encauzamiento
2016--2017}». Firman el gobernador Juan Manuel Urtubey y los intendentes de \textbf{veintiún
municipios}, entre ellos \textbf{La Caldera y Vaqueros}.

\textbf{Antecedente:} la Ley provincial 7.931 autorizó operaciones de crédito público por
hasta \textbf{trescientos cincuenta millones de dólares}.

\textbf{Cláusula primera.} El Ministerio de Ambiente por la Secretaría de Recursos Hídricos,
el Ministerio de Hacienda por la Secretaría de Asuntos Municipales y los municipios
«encararán acciones en forma conjunta a fin de dar cumplimiento a los trabajos de
encauzamiento, defensas y protección contra los desbordes de los ríos».

\textbf{Cláusula segunda.} Los proyectos «cuentan con la intervención técnica de la Secretaría
de Recursos Hídricos».

\textbf{Cláusula tercera.} Anticipo de fondos de hasta el \textbf{treinta por ciento}; el resto
contra certificados de avance conformados por Recursos Hídricos. \textbf{Los fondos se
transfieren a cuentas bancarias de los municipios.}

\textbf{Cláusula cuarta.} Los municipios se obligan a colocar un cartel de obra con el logo
del Plan del Bicentenario.
\end{ficha}

\begin{cautela}
\textbf{El decreto está publicado aparte del acta.} La ficha web del Boletín (Nº 19.913,
29/11/2016) trae el visto, los considerandos y el articulado; el artículo 1º aprueba el Acta
Acuerdo como anexo, y refrendan Gomeza, Montero y Simón Padrós. El archivo de anexos arranca
en el folio 3: los dos primeros son, por posición, el decreto en el expediente.
\textbf{Ninguna de las dos piezas transcribe la fecha del decreto ni el número de expediente.}

Faltan también dos piezas que el propio Acta invoca. \textbf{Los «proyectos de los trabajos»}
de la cláusula segunda se mencionan como preexistentes y con intervención técnica de Recursos
Hídricos, pero no se adjuntan ni se identifican por número. Y \textbf{la distribución de fondos
por municipio no consta}, ni el Acta declara dónde está.

Ninguna de las tres se reconstruye por inferencia. \textbf{El Acta dice «encauzamiento,
defensas y protección contra los desbordes de los ríos» y no nombra un río: eso es todo lo que
se sabe.}
\end{cautela}

Cuatro cosas, y tres tocan capítulos distintos de este libro.

\textbf{Primera: el municipio ejecuta y cobra.} No es un programa que la Provincia lleve
adelante sobre el territorio municipal: los fondos van a cuentas del municipio, que ejecuta
los trabajos y certifica avance. \textbf{Es la misma figura que el capítulo \ref{cap:hacienda} documenta en
1933, 2014 y 2017 con Vialidad}, ahora aplicada a los ríos.

\textbf{Segunda: se financia con deuda pública.} El antecedente invocado es una ley que
autoriza a tomar hasta trescientos cincuenta millones de dólares. Las defensas de rama y
piedra que este libro sigue desde 1910 se pagan, en 2016, con crédito internacional.

\textbf{Tercera, y es la que más importa: explica quién hace los trabajos que el Plan de
Manejo objeta.} El diagnóstico de 2024 concluye que hay «deficientes trabajos de
encauzamiento de ríos y arroyos \ldots\ realizados por el municipio y por privados sin
autorización de la Autoridad de aplicación». \textbf{El Plan de Mínimas es el marco por el
cual el municipio encauza ríos con fondos provinciales.} Que la cláusula segunda exija
intervención técnica de Recursos Hídricos y que ocho años después un informe encargado por esa
misma Secretaría diga que los trabajos se hacen sin autorización es una contradicción que este
libro registra sin poder explicar.

\textbf{Y cuarta: la única obligación de publicidad es el cartel.} El acta exige colocar un
cartel con el logo del Plan del Bicentenario. No exige publicar el proyecto, ni el monto por
municipio, ni los certificados de avance. \textbf{Lo que el ciudadano puede ver de esta obra
es el logo.}

\subsection*{Y en 2020 los fondos no llegaron}

El Plan tiene una tercera fecha, y la aporta el propio municipio.

\begin{ficha}{Municipalidad de La Caldera, 22 de octubre de 2020}
«Pese a que aún el Gobierno provincial \textbf{no habilitó los fondos del Plan de Mínima para
Defensa y Encauzamiento de Ríos}, la Municipalidad avanza con \textbf{recursos propios} en
obras de contención y encauzamiento de los arroyos \textbf{El Durazno y Guaranguay}. La empresa
\textbf{Supercemento} colabora con maquinarias.»

El intendente \textbf{Diego Ramiro Sumbay}\index{Sumbay, Diego Ramiro} informa que personal de
Obras Públicas relevó previamente los puntos críticos de desborde, que se releva el estado de
acequias y canales, y que \textbf{los sedimentos extraídos de los arroyos se usan para
acondicionar las calles} más deterioradas.

Consigna además que el municipio \textbf{gestiona los fondos desde el 31 de julio} ante la
Secretaría de Recursos Hídricos «lo que hasta el momento no se ha efectivizado», y anuncia la
reunión del \textbf{Comité Operativo de Emergencia Municipal}.
\end{ficha}

Cuatro cosas.

\textbf{Primera: el Plan de Mínimas existía y no se financiaba.} Cuatro años después del Acta
Acuerdo firmada por veintiún intendentes, el municipio dice públicamente que los fondos no se
habilitaron y que trabaja con recursos propios. Eso completa la secuencia: 2010 en ejecución,
2016 formalizado por decreto, \textbf{2020 sin fondos}.

\textbf{Segunda: los arroyos son los mismos.} El Durazno y Guaranguay, es decir los dos cursos
cuya línea de ribera se determinó en diciembre de 2014 sin publicar coordenadas ---capítulo
\ref{cap:ribera}--- y los que atraviesan el loteo del capítulo \ref{cap:loteo}.

\textbf{Tercera: quién presta las máquinas.} \textbf{Supercemento} es la contratista del
Acueducto Norte, la misma que en agosto de 2017 se reunía con la Secretaría de Recursos
Hídricos para coordinar la determinación de línea de ribera en la zona de la planta de
tratamiento. En 2020 aporta maquinaria a un municipio para encauzar dos arroyos.

\begin{cautela}
Este libro \textbf{no afirma} que haya en eso irregularidad alguna. Una empresa que tiene
maquinaria en la zona por una obra provincial y la presta a un municipio que no consiguió
fondos es una situación corriente y, en apariencia, benéfica.

Lo que registra es que \textbf{no consta ningún instrumento}: ni convenio, ni acto que acepte
la colaboración, ni constancia de intervención técnica de Recursos Hídricos, que la cláusula
segunda del Acta Acuerdo exige --- y que la reglamentación del Código de Aguas exige por su
cuenta. El artículo 236 del Decreto 2.299/03 dispone que «\textit{requerirán autorización previa
de la Autoridad de Aplicación los desmontes, \textbf{las obras de canalización, encauzamiento},
saneamientos y desagües pluviales, y toda otra que afecte o modifique el libre escurrimiento de
las aguas}». Y que este mismo capítulo documenta, cuatro años después, que el diagnóstico oficial objeta trabajos de encauzamiento «realizados por el
municipio y por privados sin autorización de la Autoridad de aplicación».

\textbf{Esta nota de prensa es la única descripción concreta que este relevamiento tiene de uno
de esos trabajos.}
\end{cautela}

\textbf{Y cuarta: el destino del sedimento.} El material extraído del cauce se usa para
acondicionar calles. Es razonable y probablemente inevitable, y al mismo tiempo es
\textbf{extracción de áridos de un cauce}, hecha por el municipio, sin que conste concesión ni
evaluación. El capítulo \ref{cap:resistencias} documenta veinte concesiones para hacer
exactamente eso, con expediente, canon e informe de impacto ambiental.

 Volviendo al acta, las otras cuatro cosas que revela son éstas.

\textbf{Las obras que sí se hicieron fueron de emergencia y no las del Plan.} El Dr.\ Le Favi
expresa que se realizaron trabajos de campo «con especial observancia del espíritu técnico del
Plan de Manejo \textbf{si bien no son idénticas a las allí previstas}». El Ing.\ Ramos precisa
que se trabajó en los barrios afectados y «en el tercio medio del río».

\textbf{El monitoreo del río lo hace una universidad privada.} «Se monitorea la tendencia
hidrogeológica de los cauces a través de una \textbf{estación de UCASAL}», y el monitoreo
aguas arriba «está previsto» en el Plan pero no realizado. El capítulo \ref{cap:defensas} documenta que en
1980 había un observatorio pluviométrico provincial con observador pago en Los Yacones.
\textbf{Cuarenta y seis años después, el dato del río lo produce una estación universitaria.}

\textbf{Los bomberos no alcanzan.} El Dr.\ Chávez afirma que La Caldera no tiene bomberos y que hubo que solicitar
la intervención de Defensa Civil. La afirmación era imprecisa ---el acta del Comité de Emergencia de marzo de 2024 registra bomberos voluntarios sin movilidad, como se ve más abajo---, y un mes después de la audiencia resultó más cierta que imprecisa: el 10 de marzo de 2026 la Inspección General de Personas Jurídicas, con Defensa Civil, encontró la sede de la Asociación de Bomberos Voluntarios Papa Francisco «cerrada con candado, sin ningún personal de guardia, y un estado de total abandono y deterioro»; un informe de Defensa Civil concluyó que la institución «no cumple con los requisitos mínimos operativos» y que eso representa «un riesgo significativo para la comunidad»; y el 12 de marzo la asociación fue intervenida, con autoridades vencidas y balances sin aprobar (Resolución 26/26, B.O.\ Nº 22.162). Describía, en todo caso, lo mismo: un municipio de cuatro mil quinientos habitantes, en un departamento de doce mil trescientos, sobre un
río que se desborda, sin un cuerpo de bomberos equipado.

\textbf{Y falta la alerta temprana, que cuesta poco.} El Dr.\ Lazarte manifiesta haber
consultado al Servicio Meteorológico Nacional y que los dispositivos de monitoreo
\textbf{«requieren de inversión mínima y permitirían prevenir nuevos eventos similares»}. Su
observación final es la que cierra la audiencia: durante la primera inundación «se adoptaron
acciones de mitigación cuando los sucesos ya estaban en curso».

\subsection*{Y dónde quedó el Plan}

El punto IV de la resolución merece leerse despacio, porque define cómo accede al documento
la gente a la que protege.

El Plan de Manejo se pone a disposición en el domicilio constituido por el colectivo actor, sito en «Avenida General Güemes S/N, \textbf{Medidor 994}» de
la localidad de La Caldera, «para que el mismo sea consultado junto con esta resolución por el
colectivo alcanzado». Y \textbf{la gestión queda a cargo de la parte actora}.

De modo que el documento que ordena cómo se maneja la cuenca que abastece a la ciudad de
Salta es un archivo digital guardado en un domicilio \textbf{cuya dirección se
identifica por el número de medidor eléctrico}, porque la calle no tiene numeración. Y quien
debe hacerlo circular no es el Estado que lo redactó ni el juzgado que lo aprobó: son los
propios vecinos que demandaron.

Este libro sostiene, y el capítulo \ref{cap:opacidad} desarrolla, que la opacidad de este dispositivo casi nunca
consiste en negar un documento sino en entregarlo de un modo que lo vuelve inutilizable.
\textbf{No hay mejor ejemplo que éste, y no lo produjo la administración: lo produjo una
sentencia que quiso ser accesible.}

\section{Qué resolvió la jueza el 12 de febrero de 2026}

El acta de esa audiencia permite decir, por fin, qué se aprobó y con qué condiciones.

\begin{ficha}{Acta de audiencia y resolución --- Expte.\ MIN 780280/22, 12 de febrero de 2026, 10:15}
Juzgado de Minas de Salta. Jueza \textbf{María Victoria Mosmann}; secretario Dr.\ Daniel
Arturo Isa.

\textbf{Intervinientes:} por la parte actora, el Dr.\ Eduardo José Lazarte Vigabriel con
patrocinio del Dr.\ Julio Chávez; por la Provincia, la Dra.\ Paula Herrando (Fiscalía de
Estado); por la Subsecretaría de Recursos Hídricos, el Dr.\ David Le Favi y el Ing.\ Ricardo
Ramos; por el Ministerio Público, el Dr.\ Agustín Vidal, Fiscal Civil Nº 2.

\textbf{Resuelve:}

\textbf{1.} Aprobar el \textbf{Plan de Manejo Integral de la Microcuenca del Río La Caldera}.

\textbf{2.} \textbf{Intimar a la Municipalidad de La Caldera} a dar cumplimiento estricto y
oportuno al \textbf{Protocolo de Crecidas e Inundaciones} por medio del COE, \textbf{bajo
apercibimiento de remitir antecedentes a la justicia penal}.

\textbf{3.} La Provincia deberá presentar en \textbf{cuarenta y cinco días corridos} la
planificación de obras y medidas prioritarias, y en un plazo posterior de otros cuarenta y
cinco el cronograma del resto de las obras previstas.

\textbf{4.} Poner el archivo digital del Plan \textbf{a disposición del colectivo afectado},
en el domicilio fijado en La Caldera.

\textbf{5.} Remitir nota de reconocimiento al \textbf{COPAIPA} y a las \textbf{veedoras} por la
labor técnica realizada.
\end{ficha}

Cuatro observaciones.

\textbf{El apercibimiento penal recae sobre el municipio, no sobre la Provincia.} Y no por
incumplir el Plan, sino por incumplir el \textbf{Protocolo de Crecidas} a través del Comité de
Emergencia. Es la obligación operativa y cotidiana, no la de planificación.

\textbf{Los plazos corren sobre la Provincia, y son cortos.} Cuarenta y cinco días para las
obras prioritarias y otros cuarenta y cinco para el resto. Contados desde el 12 de febrero,
el primero vencía a fines de marzo de 2026 y el segundo a mediados de mayo. \textbf{El
expediente registra un pedido de prórroga en mayo, concedido, y tenido por cumplido el 1º de
junio.}

\textbf{La publicidad del Plan tiene una forma precisa: un archivo digital en un domicilio de
La Caldera.} No se ordena publicarlo en el Boletín Oficial ni en el sitio del organismo. Para
acceder al documento que ordena cómo se maneja la cuenca hay que ir hasta el pueblo.

\textbf{Y el reconocimiento va al COPAIPA y a las veedoras.} El Consejo Profesional de
Agrimensores, Ingenieros y Profesiones Afines aparece así como el cuerpo técnico que acompañó
el proceso. Es la misma entidad que el capítulo \ref{cap:pdua} encuentra en el Consejo de Control del PDES 2030, y cuyas veedoras observaron el diseño del taller de 2023; integra además el Consejo Asesor del ordenamiento de bosques, y una fundación vinculada a él participó del taller provincial de ese mismo mes.

\section{Antes del Plan: lo que el expediente acumuló entre 2022 y 2023}

Cinco documentos del expediente, anteriores al Plan de Manejo, muestran cómo llegó cada organismo a la etapa de su elaboración: qué ofreció la autoridad hídrica, qué dijo la empresa de agua, cómo se diseñó la participación y con qué presentación se anunció el estudio. Ninguno es el Plan. Todos ayudan a leerlo.

\subsection*{La autoridad hídrica: un plan de monitoreo y una advertencia sobre la escala}

\begin{ficha}{Nota de la Secretaría de Recursos Hídricos del 20 de enero de 2023 (expediente 0010007-67527/2019-19, fojas 15 y 16)}
Dirigida al secretario de Recursos Hídricos, Ing.\ \textbf{Mauricio Romero Leal}\index{Romero Leal, Mauricio}; firman el Dr.\ \textbf{David Le Favi}, director de Fiscalización y Control, y la analista química \textbf{Mónica P.\ Rodríguez}. Responde a un requerimiento de Fiscalía de Estado en el amparo \textit{Lazarte Vigabriel}.

Acompaña el \textbf{Plan de Monitoreo de la Calidad del Agua Superficial de la Subcuenca Mojotoro}, parte de la red de la cuenca del Bermejo, con \textbf{siete puntos de muestreo en la microcuenca del río La Caldera}.

Y agrega dos afirmaciones de fondo. Que el Plan de Manejo Integral «debe ser \textbf{estrictamente ajustado a la escala predial} de dicha microcuenca, en función de lo establecido por la Ley Nº 7543» de bosques nativos, por lo que el plan de monitoreo «excede ampliamente estos preceptos» y se agrega «sin consentimiento de lo manifestado por los amparistas». Y que los estándares de calidad del agua para consumo humano «\textbf{son responsabilidad de los usuarios}», como resultado de la regularización de las concesiones y «de la eventual naturaleza de \textbf{sub-prestadores de servicios públicos domiciliarios} que adquirirán los responsables de las urbanizaciones».
\end{ficha}

\begin{ficha}{Plan de Monitoreo de Aguas Superficiales 2022 --- Subcuenca Mojotoro, agosto de 2022}
Secretaría de Recursos Hídricos, Dirección General de Gestión Integrada de Cuencas; Programa Calidad del Agua y Dirección de Fiscalización y Control. Continúa un \textbf{plan de rutina diseñado en 2012}: un muestreo en estiaje, entre agosto y noviembre, y otro en época de lluvias.

Puntos en la microcuenca del río La Caldera: río Santa Rufina aguas abajo de la confluencia con el San Alejo, en la toma del dique Campo Alegre (ASMo~2); salida del dique (ASMo~3); río La Caldera aguas arriba del camping (ASMo~4) y dos kilómetros aguas abajo del puente de ingreso a la localidad (ASMo~5); río Wierna antes de la confluencia (ASMo~6); río La Caldera en el puente Wierna (ASMo~7); y río Castellanos aguas abajo de la confluencia con el Lesser (ASMo~8). Aguas abajo siguen el río Vaqueros, el Mojotoro, el desagüe de la zona norte de Salta, la planta depuradora norte, Campo Santo, la Ruta 34 y la descarga de la central térmica de Güemes.

Parámetros: temperatura, conductividad, pH, oxígeno disuelto, cloruro, amonio, nitrito, nitrato, fósforo, DBO$_5$ y coliformes totales y fecales. El documento anuncia un informe final de campaña \pendiente{resultados}.
\end{ficha}

\begin{cautela}
\textbf{Dos lecturas de la misma nota.} La primera: la autoridad hídrica mide desde 2012 la calidad del agua en la cuenca con una red propia, y este libro no había visto esa red; sus resultados son el dato que falta para discutir con números lo que el capítulo \ref{cap:aguabaja} discute con testimonios.

La segunda es más delicada. Leer la Ley 7543 como un límite a la «escala predial» es una interpretación del organismo, no una cita de la ley, y define el alcance del Plan antes de que se escriba: \textbf{un plan de cuenca reducido a la escala del predio}. Y la frase sobre los urbanizadores anticipa, en 2023 y por escrito, lo que el capítulo \ref{cap:aguabaja} documenta como práctica: que el agua de los loteos quede a cargo de sus responsables, convertidos en prestadores.
\end{cautela}

\subsection*{La empresa de agua: dos acueductos, veintidós mil usuarios y un río que no alcanza}

\begin{ficha}{Informe de Aguas del Norte (CoSAySa) para Fiscalía de Estado, con plano del 2 de diciembre de 2022}
Firma el Ing.\ Arnaldo Alfaro, jefe de Producción, junto con la jefatura del Departamento Servicio Capital.

\textbf{Sistemas sobre el río La Caldera:}
\begin{itemize}[leftmargin=*,nosep]
\item \textbf{Acueducto Norte}: captación en la margen derecha del río La Caldera, a la altura de la curva de El Gallinato; \textbf{1.200 metros cúbicos por hora} y \textbf{10.822 usuarios} en barrios de la capital (15 de Setiembre, Juan Manuel de Rosas, Leopoldo Lugones, Ciudad del Milagro, Castañares, Parque Belgrano II y III, cisterna Huaico).
\item \textbf{Acueducto Wierna}: captación en las márgenes activas del río Las Nieves (drenes de la UNSa y «Manantiales»); \textbf{900 metros cúbicos por hora} y \textbf{11.587 usuarios}, también en la capital.
\item \textbf{Acueducto Campo Alegre}: desde el dique; a la fecha del informe faltaban «detalles menores» para que operara.
\end{itemize}

\textbf{Problemas que declara la empresa:} para sostener 1.200 metros cúbicos por hora el río debe llevar más de 3.500; en estiaje se pide al Consorcio del Río Mojotoro que abra la válvula de chorro hueco del dique, y el consorcio prefiere hacerlo a mediados de agosto; hay cámaras dentro de matrículas privadas y propietarios que no dejan pasar el agua para alimentar los drenes; los cambios de cauce erosionan y dejan cañerías al descubierto; y «se están llevando a cabo \textbf{urbanizaciones privadas} los cuales no tienen un sistema de tratamiento de efluente cloacal por lo que \textbf{pueden llegar a contaminar las napas}». En el Wierna, aunque las instalaciones están dentro de las líneas de ribera, vecinos de matrículas adyacentes «alambraron \ldots\ invocando derecho de propiedad».

\textbf{Antecedente:} a fines de 2021, un acuerdo entre la Secretaría de Minería y CoSAySa (expediente 99078-2021) fijó una \textbf{nueva zona de exclusión} para canteras.

\textbf{Plano:} «Relevamiento Acueducto Norte --- Líneas de ribera, cámaras acueducto y canteras», escala 1:20.000. Entre sus referencias figura una «Línea Ribera --- Res.\ 236--360 SRH».
\end{ficha}

\begin{cautela}
\textbf{Esto pone número a una de las tesis del libro.} El agua que se capta en el departamento abastece a más de veintidós mil usuarios de la capital; ninguno de los barrios que el informe enumera está en La Caldera. Y la empresa que la capta describe, en su propio informe, los tres riesgos que este libro sigue en otros capítulos: la extracción de áridos, el alambrado de las márgenes y los loteos sin tratamiento cloacal.

Tres datos piden cotejo. El primero: el informe ubica la captación del Acueducto Norte en el río y consigna que el Acueducto Campo Alegre todavía no operaba, mientras que la planta del dique se había inaugurado oficialmente en 2021 bajo el nombre de «Acueducto Norte»; el capítulo \ref{cap:expropiacion} pone las dos fuentes juntas. El segundo: el Acueducto Norte se anunció en 2017 con 1.100 metros cúbicos por hora (capítulo \ref{cap:ribera}); en 2022 la empresa declara 1.200. Y el tercero: la «Res.\ 236» del plano podría ser la Resolución 236/14 de la Secretaría, que rectificó la 142/14 sobre los ríos Mojotoro y La Caldera, una de las dos de la serie cuyas coordenadas no se publicaron (capítulo \ref{cap:ribera}); el plano no lo dice, y el apéndice \ref{cap:pedidos} lo pregunta.
\end{cautela}

\subsection*{La participación, tal como se diseñó}

\begin{ficha}{Diseño de Taller Participativo, 31 de agosto de 2023}
Previsto para el \textbf{29 de septiembre de 2023}, de 8:30 a 14:00, en el Centro Gaucho El Palenque (capacidad, trescientas personas), «como instancia previa al diseño y presentación del Plan de Manejo Integral». Invoca las sentencias del \textbf{17 de diciembre de 2021} y del \textbf{7 de julio de 2023}.

\textit{Esa segunda sentencia aparece aquí y no vuelve a aparecer. Es la única mención que el corpus de este libro tiene de una decisión judicial posterior a la condena y anterior al taller ---dieciocho meses después de la primera y ocho semanas antes del diseño que la invoca---, y no se sabe qué resolvió: si fijó plazos, si resolvió una apelación, si precisó el alcance de las dos etapas que las partes discutirían en 2026. \textbf{Es, después del texto íntegro de la condena, el documento más importante que falta de esta causa}, y el apéndice \ref{cap:pedidos} lo pide.}

\textbf{Mesa informativa:} exposición del juzgado; «línea de base» y «elaboración del Plan», a cargo de las secretarías de Recursos Hídricos, de Minería y Energía y de Ambiente, del municipio y de la consultora (Ing.\ Javier Ramos). Luego, trabajo grupal con la técnica del árbol de problemas, plenario y formulario individual.

\textbf{Convocados, por sector:} los actores del amparo; vecinos de la zona de influencia; los \textbf{loteadores} del municipio ---club de campo «Potrero de La Caldera» (Ramiro Costas de Arguibel), «Urbanización Valle Alegre» (Luis Cabanilla), «La Milagrosa» (Matías Marcelo Cornejo Isasmendi), «El Durazno» (\textbf{Renzo Palacios} ---el primer socio gerente de Jardín Celestial S.R.L.\ y titular del 24,99\,\% de sus cuotas desde 2014, capítulo \ref{cap:loteo}---), «Ayeres de Quitilipi» ---así en el diseño; «Ayres» en las notas municipales--- (Augusto Torino), «Vistas del Lago» (Federico Antonio Fernández) y un loteo sin nombre en La Calderilla (Juan Geracaris)---, además del Barrio 20 Viviendas; pobladores rurales de los parajes Potrero de Castilla, Los Yacones, Wierna, San Alejo, Santa Rufina y Cuchimayo, y la escuela 4534 de Los Yacones; tres \textbf{productores de áridos} (la cooperativa La Mesa Redonda, Alejandrina Prudencia Morales y Geomix S.R.L.); las comunidades \textbf{Kolla Kondorwaira} y \textbf{Kolla Urkuwasi}; y el Consorcio de Usuarios del Río Wierna.

\textbf{Compromisos de registro:} desgrabación a cargo del municipio, sistematización por observadores de Recursos Hídricos y de Minería, informe final al juzgado y publicación como devolución a la comunidad.
\end{ficha}

\begin{cautela}
\textbf{La lista de convocados es, en sí misma, un mapa de actores.} Por primera vez un documento oficial reúne en una página a los loteadores del municipio con nombre y emprendimiento, a los productores de áridos y a las dos comunidades indígenas. El Durazno figura entre los loteos, y su loteador entre los invitados.

El taller se postergó. Una segunda versión del diseño, fechada en septiembre, lo pasa al \textbf{10 de octubre de 2023}, y así lo comunica la invitación del 25 de septiembre a los actores del amparo. Esa versión fija también los plazos que faltaban: desgrabación y sistematización hasta el 9 de noviembre, e \textbf{informe final al Juzgado de Minas hasta el 27 de diciembre de 2023}, lo que concuerda con el informe de diciembre que registra el expediente judicial. El informe final, con su propia fecha de 2023, se obtuvo y se analiza más abajo.
\end{cautela}

\subsection*{Cómo se preparó el taller}

\begin{ficha}{Expedientes administrativos y judicial, agosto a octubre de 2023}
\textbf{7 y 17 de agosto.} La Municipalidad remite a Fiscalía de Estado (nota 0010007-22011/2023-0) los convocados que le corresponde aportar: los siete loteadores; un relevamiento de su Subsecretaría de Asuntos Rurales, titulado «riesgo de inundaciones de ríos y arroyos en la época estival», con los pequeños productores de los parajes; y dos funcionarios municipales. \textbf{La lista de loteadores del diseño es la que dio el municipio.}

\textbf{7 de septiembre.} Las veedoras, representadas por el COPAIPA (Ing.\ María Lorena Portal), elevan al juzgado doce observaciones al diseño: un formulario de inscripción en papel en la Municipalidad, los responsables de la difusión, la normativa sobre pueblos originarios, facilitadores para el trabajo grupal, plazos para el informe final y «disponer del contrato de la empresa consultora a cargo del Plan, para el momento del taller».

\textbf{18 de septiembre.} Un decreto de la jueza Mosmann (actuación 9.919.447) intima a la Provincia a responder en tres días las observaciones de las veedoras (fs.\ 266); a informar, en el mismo plazo y a pedido del actor (fs.\ 275/276), «el procedimiento y/o criterio empleado para la selección de actores a ser convocados al taller participativo»; y, sobre el acta de la reunión con las comunidades que la Provincia había acompañado cinco días antes (fs.\ 278/283), resuelve \textbf{no proveer} hasta que acompañe «constancia de la consulta efectuada a la Secretaría de Asuntos Indígenas» ---la que la propia Provincia había anunciado en el punto 7 de su informe de fs.\ 228/231--- y acredite «la calidad de autoridades y/o representantes de las comunidades involucradas respecto de los participantes de dicha reunión». Es decir: el acta sola no le alcanzó al juzgado ---una semana, como se verá---.

El informe de fs.\ 228/231 (actuación 9.712.229), que la Provincia había presentado después de la audiencia del 2 de agosto, es el origen de la lista de convocados que el diseño del taller reproduce: ocho sectores ---los actores, los vecinos, veinte usuarios de agua con nombre entre los que figuran dos granjas avícolas y María Eugenia Carante (capítulo \ref{cap:redes}), los loteadores, los productores rurales relevados por el municipio, tres canteras, las dos comunidades y Aguas del Norte «en virtud de la zona de exclusión»---. \textbf{Enumera, pero no dice con qué criterio se eligió a cada uno}, y se reserva ampliar la lista «a consideración de esta parte». Es exactamente lo que el actor pidió precisar y lo que el decreto del 18 de septiembre intimó a informar.

La respuesta llegó el \textbf{21 de septiembre} («Acompaño Metodología --- Aclaraciones», actuación 9.946.997), con la metodología corregida según las observaciones de las veedoras. Sobre la selección, dice que cada organismo ---Ambiente, Minería, Recursos Hídricos, la Municipalidad y Asuntos Indígenas--- informó los sectores con implicancias en la zona «de acuerdo a los registros de cada organismo», y suma a la lista las Cámaras Legislativas, el Concejo Deliberante y el intendente; agrega que la difusión alcanzará «a toda persona que tenga algún interés». \textbf{Es un criterio de procedencia, no de selección}: dice de qué registro salió cada nombre, no por qué entraron unos y no otros. Sobre la consultora, el mismo escrito informa que «oportunamente se informará la empresa consultora que estará a cargo de la elaboración del Plan de Manejo Integral o bien, el equipo de técnicos» de las secretarías y el municipio: \textbf{tres semanas antes del taller, la Provincia todavía no le decía al juzgado quién iba a hacer el Plan}.

El \textbf{25 de septiembre} (actuación 9.972.217) el juzgado da todo por cumplido. Aprueba la metodología definitiva, con un cronograma que declara «tentativo», y registra que la actora no había objetado la propuesta original. Sobre la selección de actores, «téngase presente lo manifestado»: no la observa. Y sobre la consulta, da «por cumplida la intimación de fs.\ 284, párrafo sexto» ---la que exigía la constancia de Asuntos Indígenas \textit{y} la acreditación de la calidad de autoridades o representantes de quienes firmaron por las comunidades--- con la documentación de fs.\ 317/333 (actuación 9.947.016). Para la representatividad, lo que la Provincia había ofrecido era el «listado de las autoridades que participaron de la mencionada reunión», y lo que la actuación 9.947.016 contiene son las dos invitaciones, la planilla y el acta, sin el acta de asamblea que el dictamen del 4 de septiembre había recomendado. \textbf{Una semana después de exigirla, la acreditación se tuvo por hecha con la nómina de asistentes.}

\textbf{20 de septiembre.} El Dictamen 129/2023 de la Dirección de Fiscalización y Control de Recursos Hídricos (Dr.\ David Le Favi; expediente 0090034-214840/2023-0) responde por el sector hídrico: en el departamento hay \textbf{275 usuarios de agua pública para riego}; como no todos están dentro del área de la cuenca, la Secretaría estaba determinando, con las herramientas de la Dirección General de Inmuebles, qué catastros con derecho de riego caen en la zona, para notificarlos. Acompaña el padrón de regantes del departamento.

\textbf{La versión de septiembre del diseño} incorpora casi todas las observaciones y suma a los «usuarios y propietarios colindantes al dique Campo Alegre». La exposición central pasa a llamarse «Fase Preliminar de la elaboración del Plan de Manejo Integral», y el nombre del consultor sale de la agenda. En la invitación, «medidas estructurales y no estructurales para mitigar el impacto de las lluvias y desbordes fluviales» se convierte en «medidas integrales».

\textbf{Sin fecha legible.} El área de Acción Social municipal relevó a \textbf{veintiocho familias} con fichas sobre su situación habitacional ---si seguían en su casa, la habían abandonado, estaban en un albergue temporal o con parientes---. Contienen datos personales y no se transcriben.
\end{ficha}

\begin{cautela}
\textbf{Tres observaciones.} La primera: para el sector hídrico, el criterio de convocatoria fue el padrón de riego. El agua de la cuenca quedó representada por quienes tienen concesión para usarla, no por quienes sufren sus crecidas; a éstos los aportó el municipio, por su propio relevamiento.

La segunda: las veedoras pidieron disponer del contrato de la consultora para el momento del taller, y el consultor salió de la agenda de la versión de septiembre. La contratación se publicó recién el 26 de enero de 2024 (B.O.\ Nº 21.640; véase la cronología). Y sin embargo, el 10 de octubre de 2023 el consultor expuso en el taller, se presentó como «contratado por la provincia» y fijó el área del estudio por «los alcances de la contratación»: \textbf{tres meses antes de que esa contratación se publicara} \pendiente{fecha de firma}.

La tercera: la versión que llegó a los vecinos ya no nombra las lluvias ni los desbordes. Es un cambio de redacción, pero es el que convierte un plan contra las inundaciones en un plan «integral».
\end{cautela}

\subsection*{Lo que pasó en el taller}

\begin{ficha}{«Taller participativo La Caldera. Informe final», 2023 (61 páginas)}
\textbf{Participación.} El taller se hizo el 10 de octubre de 2023 en el Centro Gaucho El Palenque. Hubo 165 inscripciones y \textbf{88 personas acreditadas}: 22 por los sectores convocados (loteos, comité de emergencia municipal, canteras, comunidades originarias y productores) y 66 vecinos. El espacio virtual de consultas previas estuvo abierto once días y no recibió ninguna. Se completaron 48 encuestas. Nadie dijo haberse enterado por la radio, y la mitad consideró inadecuado el horario de un día hábil por la mañana.

\textbf{El marco.} La jueza María Victoria Mosmann explicó que el proceso nació de las inundaciones de diciembre de 2018 y enero de 2019 en El Durazno, El Nogalar y Potreros de La Caldera, y que el taller no era una audiencia pública.

\textbf{El estudio.} El consultor, Ing.\ Javier Ramos Vernieri, de YSATÍ, lo presentó como el «Plan director de desagües pluviales y fluviales de la localidad de La Caldera», con un plazo de ciento veinte días. Su área de intervención es la zona urbana de la margen derecha. A la pregunta por La Calderilla y la margen izquierda respondió que la ampliación era factible, pero que «no está dentro de los alcances de la contratación». Fiscalía de Estado agregó que el área era la que fijaba la sentencia, y el coordinador de Recursos Hídricos, David Le Favi, que se tomaría nota. Una vecina de La Calderilla contó que el año anterior el río se había llevado más de la mitad del terreno de su padre, y que con la extracción de áridos «el río se está tirando para el lado de La Calderilla».

\textbf{Los pedidos.} Varios vecinos pidieron que, mientras se elaborara el plan, se suspendieran la extracción de áridos, los desmontes y la venta de lotes, con el aplauso de la sala. Fiscalía respondió que la suspensión no estaba contemplada. El abogado de los amparistas recordó las tres acciones que el amparo identificó como las más degradantes: la extracción de áridos, la falta de control inmobiliario y la obra de Supercemento, de la que ---dijo--- «hay constancia en el expediente» de que desmontó bosque ribereño. Un vecino planteó la participación de funcionarios o de sus familiares en las empresas loteadoras y extractoras: «si la misma persona está de los dos lados del mostrador, el medio ambiente no va a mejorar».

\textbf{Los grupos.} Se formaron cuatro. \textbf{El de los funcionarios se retiró antes del final y no elaboró su árbol de problemas.} Los otros tres coincidieron en las causas: falta de control y de control cruzado entre organismos; extracción de áridos fuera del tercio medio, «donde decide el maquinista»; loteos sin estudio de impacto; certificados de no inundabilidad sin manejo del agua; catastros otorgados sin esa certificación; municipios que aceptan la donación de calles y convierten el loteo en barrio abierto; líneas de ribera sin marcar y propietarios que «afirman que son dueños de la ribera»; y pozos ciegos que contaminan el acuífero. Propusieron una comisión de vecinos y municipio para controlar la extracción, y una mesa de diálogo, para la que dejaron sus nombres.

\textbf{Las encuestas.} En orden de frecuencia, las causas fueron la explotación minera, los desmontes, la falta de control, los loteos y el incumplimiento de la ley. Hubo una mención al acueducto Norte: «se llevan recursos para la ciudad y dejan a los pobladores sin agua».
\end{ficha}

\begin{cautela}
\textbf{El informe registra la exclusión y la desmiente en la misma página.} Su apartado de devolución justifica la reducción del área de estudio porque la mayoría de las «necesidades hídricas registradas» está en la margen derecha urbana, y agrega que «tampoco en el marco del proceso se han recepcionado planteos específicos de esta área». Pero el mismo informe transcribe, páginas antes, los planteos de los vecinos de La Calderilla y de la margen izquierda, y las encuestas piden sumar La Calderilla, Los Yacones y el Nevado de Castilla. Y el \textbf{acta del taller}, firmada ese mismo 10 de octubre por la moderadora y los funcionarios presentes ---se abrió a las 9.21 y cerró a las 13---, dejó escrito que el pedido de ampliar lo resuelto «a la margen izquierda del Río» \textbf{«será volcado en el Informe Final»}. El acta registra también que los vecinos pidieron a las partes la paralización inmediata de los desmontes y de las actividades extractivas hasta la presentación del Plan, y que, sobre la obra del acueducto, \textbf{manifestaron no haber tenido instancias participativas}.

\textbf{Y la suspensión que pidió la sala no tiene destinatario.} La Secretaría de Ambiente respondió que sólo puede suspender caso por caso, dentro de un sumario, y que una medida general «deberá ser evaluada y dictada por el Organismo competente», sin decir cuál.
\end{cautela}

\subsection*{La devolución: lo que respondió cada organismo}

\begin{ficha}{Informe final del taller, apartado «Devolución a la comunidad»}
\textbf{Municipalidad.} Enumera ocho acciones: limpieza y encauzamiento de cauces; pedido anual de planes de mínima; asistencia social a damnificados; creación de la Junta de Defensa Civil y del Comité Operativo de Emergencia, con un protocolo aprobado; pedido a Minería de un convenio temporario de áridos que la obligue a controlar y denunciar; actas de clausura remitidas a Ambiente, y \textbf{suspensión del trámite de los loteos que no presenten su documentación}; pedido de audiencia al Ministerio de Obras Públicas porque los \textbf{planes de mínima de 2022 y 2023 seguían sin aprobarse}; y reunión del comité de emergencia con los vecinos.

\textbf{Secretaría de Ambiente.} Su poder de policía es «normativo y reglamentario, careciendo del ejercicio de la fuerza pública». Los proyectos de loteo se evalúan ante el municipio, con vista obligatoria a la Secretaría. En zonas de valor medio o alto del OTBN, el cambio de uso del suelo exige su certificado previo (art.\ 20 de la Ley 7543), y sobre la Ruta 9 se requiere además la autorización de Vialidad Nacional, que prohíbe toda apertura en camino de cornisa.

\textbf{Secretaría de Minería y Energía.} Informa las canteras del área y el estado de sus informes de impacto ambiental. En el río La Caldera:
\begin{itemize}[leftmargin=*,nosep]
\item Tierra Gaucha (Geomix S.R.L., expediente judicial 20.738): informe \textbf{insuficiente}; en readecuación de área por la zona de exclusión.
\item La Serena (Morales, 18.995): informe \textbf{insuficiente}; mensurada.
\item Mesa Redonda (cooperativa, 14.971): informe \textbf{insuficiente}; concesión renovada en 2016.
\item Mesa Redonda I (cooperativa, 18.134): informe aprobado.
\item \textbf{Zona de exclusión de Aguas del Norte} (Resolución SMyE 20/2022): \textbf{no presenta informe}.
\item Zona de exclusión del municipio de La Caldera (expediente administrativo 302-186719/2021): no presenta informe.
\end{itemize}
En el río Wierna figuran Hanna (provisoria), Terranostra (concesión vencida en diciembre de 2021), Los Yacones (Incovi S.R.L., renovación rechazada en junio de 2020), El Mercado 2 (AB Construcciones S.R.L., provisoria) y convenios del municipio de Vaqueros, dos de ellos ya no vigentes y dos en trámite.

\textbf{La zona de exclusión} es un área de 591 ha 2.507 m² con tomas y obras de captación. Por el convenio entre la Secretaría y CoSAySa, allí la empresa sólo realiza, por sí o por terceros, la limpieza y el mantenimiento del cauce, tres veces al año, y se compromete a denunciar cualquier extracción. La figura surge del artículo 138 de la Ley 7141, que prohíbe extraer a menos de 800 metros aguas arriba y 200 aguas abajo de una toma, salvo convenio.

\textbf{Inspecciones de 2022 y 2023:} treinta, en las que se ordenaron varias paralizaciones. \textbf{En la zona de exclusión se encontraron extractores sin permiso cuatro veces}: el 9 de marzo de 2022; el 10 de mayo de 2023, con tres extractores en el Wierna y tres en el La Caldera; y el 18 de agosto de 2023, con otros tres. Y la cuarta, el 1 de agosto de 2023, en el área de exclusión del acueducto, \textbf{«se encontraba la Municipalidad de La Caldera, se paralizaron las actividades por estar en un sector no autorizado»}. El 7 de enero de 2022, en Tierra Gaucha, operaba la \textbf{«empresa Supercemento con guías de tránsito GEOMIX»}. También se paralizaron por falta de permiso la cantera Teresa (marzo de 2022), AB Construcciones en La Serena (agosto de 2023, luego habilitada) y Hanna (julio de 2023).

\textbf{Los términos de referencia del Plan}, transcriptos al final, definen el objeto como el «Plan de Manejo de Cuencas del área urbanizada de La Caldera», con un modelo digital de terreno de 30 metros de resolución. El estudio hidrológico promete dimensionar «cunetas, alcantarillas y otras obras de drenaje necesarias para \textbf{proteger adecuadamente el camino}», una fórmula de estudio vial, y entre los organismos a consultar nombran al consorcio de riego del Valle de Siancas.
\end{ficha}

\begin{cautela}
\textbf{Cuatro lecturas.} La primera: el área donde la Provincia restringió la extracción para proteger el agua de la capital, de casi seiscientas hectáreas, está a cargo de la empresa de agua, que no presentó informe ambiental y cuyo sector fue ocupado por extractores sin permiso una y otra vez. Una de esas veces, el extractor fue el propio municipio.

La segunda: la empresa que el abogado de los amparistas señaló en el taller aparece, en la inspección de enero de 2022, operando con las guías de otra concesionaria.

La tercera: tres de las cuatro canteras que ese informe lista en el río La Caldera ---el cuadro de 2021 del capítulo \ref{cap:resistencias} registra más sobre el mismo tramo, con otro recorte--- tenían en 2023 su informe de impacto ambiental calificado como insuficiente, y seguían concedidas.

La cuarta: el acto que formalizó el acuerdo entre Minería y Aguas del Norte, que este libro buscaba, es la \textbf{Resolución 20/2022 de la Secretaría de Minería y Energía}. No figura en el Digesto Normativo Minero Ambiental de la Provincia, ni en su edición de febrero de 2025, ni entre las resoluciones número 20 de 2022 que indexa el Boletín Oficial, que corresponden a Rentas, Obras Públicas, contrataciones y personal, ni en el listado de resoluciones de la Secretaría de Minería publicadas en el Boletín. En febrero de 2022 la Secretaría de Minería y Energía estaba a cargo de la Ing.\ Flavia Royón. \textbf{Una zona de exclusión de casi seiscientas hectáreas sobre el río se rige por un acto que no fue publicado} \pendiente{texto y plano}.

\textbf{Y dos meses después del taller cambió la regla sobre esos informes.} El artículo 34 del Código de Procedimientos Mineros que el informe transcribe mandaba presentar el informe de impacto ante el Juzgado de Minas, con vista a la Secretaría de Ambiente. La \textbf{Ley 8405}, publicada el 22 de diciembre de 2023, lo reemplazó: el informe se presenta ante la Secretaría de Minería, que lo remite al organismo ambiental para su evaluación, y \textbf{las observaciones de éste no son vinculantes}, aunque su rechazo debe fundarse. Desde entonces, la misma secretaría que concede el convenio de extracción y cobra la cuota por hectárea es la que aprueba el informe ambiental de la cantera. Esa cuota fija anual fue de \$810 por hectárea en 2022, de \$1.945 en 2023 y de \$45.520 desde 2025 (Resoluciones SMyE 86/2021, 196/2022 y 134/2024, esta última del 23 de diciembre de 2024); la Resolución 176/25 de la Secretaría de Minería, del 29 de diciembre de 2025 (B.O.\ Nº 22.121 del 30/01/2026), la llevó a \textbf{\$65.320 por hectárea desde 2026}, con el índice de la FACPCE como referencia. Lo recaudado va, por disposición de esa misma resolución, al \textbf{Fondo Especial de Promoción Minera}: lo que pagan las canteras financia la promoción de la actividad. Y el manual de misiones y funciones de la Secretaría, aprobado por la Resolución 36/22 del Ministerio de Producción (B.O.\ Nº 21.178 del 22/02/2022), asigna al mismo \textbf{Programa Gestión y Policía Minera} la tarea de «crear las condiciones necesarias para el fortalecimiento de la actividad minera», la de «controlar el estricto cumplimiento de las normas mineras y ambientales» y la de evaluar los informes de impacto ambiental. Su Subprograma Áridos asesora a los municipios «sobre las posibilidades de extracción de áridos en su jurisdicción, mediante los convenios pertinentes» y, a la vez, debe velar por el cumplimiento de la norma con controles e inspecciones. \textbf{Promover, autorizar, evaluar e inspeccionar la extracción es tarea de una misma oficina.} El Digesto que recopila estas normas lo coordinaron, en sus ediciones de 2023 a 2025, la jueza de Minas que tramita este amparo y la secretaria de Minería y Energía.
\end{cautela}

\subsection*{Después del taller}

\begin{ficha}{Documentos de octubre de 2023 a marzo de 2024}
\textbf{27 de septiembre de 2023.} El Ministerio de Producción invita al taller al diputado Gustavo Javier Pantaleón y al senador Héctor Miguel Calabró, representantes del departamento.

\textbf{2 de octubre de 2023.} La Dirección General de Concesiones Hídricas eleva a Fiscalización y Control un cuadro de observaciones al padrón del Consorcio de Usuarios «Mojotoro y Regantes del Río La Caldera», para regularizarlo. Sus categorías son elocuentes: concesiones de riego sobre catastros que «no figura[n] en Inmuebles», sin número de catastro, con cambio de titularidad y, sobre todo, marcadas como \textbf{«Urbano»}, con la lista de los nuevos catastros en que se subdividió la parcela regada. El pedido al consorcio para que notificara a esos usuarios databa de marzo de 2023.

\textbf{17 de octubre de 2023.} La Secretaría de Minería y Energía remite las cédulas con que el 22 de septiembre invitó al proceso a las tres canteras (la cooperativa La Mesa Redonda, Geomix S.R.L.\ y Alejandrina Morales).

\textbf{18 de febrero de 2024.} Un informe municipal registra lluvias nocturnas que anegaron El Balcón, El Durazno y El Nogalar; el arroyo Guaranguay «sumamente caudaloso»; y el traslado de vecinos con una retroexcavadora municipal.

\textbf{3 y 4 de marzo de 2024.} Un temporal y el desborde del río dañan el loteo Santiago Apóstol; Desarrollo Social asiste a cinco familias con pérdidas materiales.
\end{ficha}

\textbf{Las concesiones de riego sobre tierras loteadas son el eslabón que faltaba} entre el padrón de agua y el negocio inmobiliario: la parcela se subdivide, pero el derecho de riego sigue inscripto como si fuera una finca \pendiente{cuadro completo de observaciones al padrón del Consorcio de Usuarios «Mojotoro y Regantes del Río La Caldera», Dirección General de Concesiones Hídricas, 2 de octubre de 2023}. El procedimiento para dividirlo existe y se publica ---el capítulo \ref{cap:aguabaja} registra casos de 2005, 2018, 2021, 2022 y 2023---; lo que no se sabe es qué parte del padrón pasó por él. Y el primer verano después del taller volvió a inundar los mismos barrios que originaron el amparo.

\subsection*{La consulta a las comunidades kolla}

\begin{ficha}{Expediente 0010007-67527/2019-28 (Fiscalía de Estado; tramitado también ante la Secretaría de Asuntos Indígenas)}
\textbf{30 de agosto de 2023.} Fiscalía de Estado (Dra.\ Paula Herrando) pide la intervención de la Secretaría de Asuntos Indígenas para la consulta previa, libre e informada, e identifica a las comunidades Kolla Kondorwaira y Kolla Urkuwasi.

\textbf{4 y 5 de septiembre.} El consultor de la Subsecretaría de Regularización Territorial dictamina que la afectación directa surge «de manera palpable y evidente»: Kondorwaira tiene relevamiento territorial por la Ley 26.160, y Urkuwasi, aunque tramita su personería, conserva el derecho a ser consultada. Recomienda que la comunicación quede «reflejada en actas comunitarias refrendadas en Asamblea comunitaria». El subsecretario, Dr.\ Francisco Ariel Sánchez, convoca a una reunión informativa; las comunidades son notificadas el 6.

\textbf{12 de septiembre, 10 horas, Municipalidad de La Caldera.} Acta manuscrita de cinco hojas ---tres de texto, una de firmas y una planilla de asistentes---. Invoca el Convenio OIT 169 y el amparo «del Expte.\ judicial 23.860/19 y 78280/22», así, con el segundo número truncado. Firman o se registran funcionarios de seis organismos provinciales: la Secretaría de Minería y Energía (Luciana Cerúsico y otra funcionaria), Recursos Hídricos (David Le Favi), Fiscalía de Estado (Paula Herrando), la Secretaría de Ambiente (Silvina B.\ Borelli, directora de Fiscalización y Control), el Ministerio de Producción y Desarrollo Sustentable y la Secretaría de Asuntos Indígenas (Francisco Ariel Sánchez); por la Municipalidad, el intendente Diego Ramiro Sumbay y su apoderado en la causa, Héctor Velásquez; y \textbf{dos representantes por cada comunidad}. Sobre la línea de base el acta dice dos cosas distintas: en la primera hoja, que «luego de la realización de talleres participativos, se presentará al Juzgado interviniente, para su aprobación»; en la tercera, que «ya se encuentra presentada al Juzgado y se encuentra aprobada». Consta también que se entrega a las comunidades «el protocolo de emergencia que ya se encuentra aprobado». Se entrega un resumen y un croquis de las zonas de impacto; se invita al taller del 10 de octubre y se deja abierto el espacio de consulta si las comunidades lo piden. Y consta: «En cuanto a las dudas que surgen respecto de \textbf{posibles obras que se realizaron posiblemente sin autorización o CPLI}, se solicitan que se realicen las consultas por escrito para poder identificar las mismas, y para que sea evaluado por el Consultor Ambiental interviniente». Las autoridades comunitarias manifiestan que participarán.

\textbf{19 de septiembre.} Un nuevo dictamen del mismo consultor tiene la consulta por cumplida: se convocó «debidamente a las autoridades comunitarias y a algunos integrantes de las Comisiones Directivas de las comunidades, conforme constancias de convocatoria (fs.\ 52/53) y asistencias (fs.\ 54)», y ningún miembro expresó «formalmente su RECHAZO». Concluye que «el objeto del presente expediente se encuentra concluido». \textbf{No se agrega ningún acta de asamblea comunitaria}: las constancias son las dos invitaciones notificadas, la planilla de asistentes y el acta de la reunión.
\end{ficha}

\begin{cautela}
\textbf{Dos precisiones sobre el acta.} La primera: el protocolo de emergencia que en septiembre de 2023 se entrega como «ya aprobado» es el mismo que en marzo de 2024 el Comité de Emergencia lleva reformado a los concejales «para aprobar por ordenanza», y cuya ordenanza de aprobación el apéndice \ref{cap:pedidos} todavía pide. La segunda: el acta cita el expediente de ejecución como «78280/22», con un dígito de menos; y no es el único número que estos papeles abrevian. El expediente administrativo, cuyo número completo es el 0010007-67527/2019-28, aparece abreviado de cinco maneras en sus propias piezas: «007-67527/2019-0» y «10007-67527/2019-28» en los dos informes de Asuntos Indígenas del 10 de agosto de 2023; «07-67527/2019-28» en la nota de Fiscalía del 30 de agosto; \textbf{«007-6752/2019-28»} ---así, con un dígito de menos--- en las invitaciones que la Subsecretaría cursó a las dos comunidades el 6 de septiembre; y «07-67527/2019» en el escrito de Fiscalía del 21 de septiembre. La forma truncada que el título de esta ficha usaba en una versión anterior no era un error de este libro: es la que llegó a las comunidades. Es uno solo, abierto en la mesa de entradas de Fiscalía de Estado el 9 de agosto (capítulo \ref{cap:resistencias}).

\textbf{Una precisión sobre la fecha}, que el capítulo \ref{cap:resistencias} también recoge: el Plan de Manejo fecha la consulta el 23 de septiembre, pero la convocatoria, las notificaciones, el acta y los dictámenes la sitúan el \textbf{12 de septiembre de 2023}.

\textbf{Y una distancia entre dos dictámenes del mismo organismo.} El del 4 de septiembre pedía actas refrendadas en asamblea comunitaria; el del 19 da por concluida la consulta con una reunión informativa y la firma de dos representantes por comunidad, cuya representatividad el juzgado pidió acreditar el 18 de septiembre y dio por acreditada el 25, con el listado de asistentes. Las comunidades, además, llevaron a esa mesa una pregunta que este libro también hace ---qué obras se hicieron sin autorización ni consulta---, y la respuesta fue que la formularan por escrito \pendiente{presentaciones de las comunidades}.
\end{cautela}

\subsection*{Lo que la Provincia y el municipio ya sabían en 2020 y 2022}

\begin{ficha}{Notas 101 a 105/20 de la División Tierra y Hábitat de la Municipalidad, octubre de 2020}
Dirigidas a los responsables del club de campo Potreros de La Caldera, de la urbanización Valle Alegre, de Vistas del Lago y de Ayres de Quitilipi, entre otros. Les piden en diez días el plano del proyecto, el certificado de no inundabilidad, los planos aprobados por Inmuebles, el informe de impacto ambiental, la factibilidad de servicios y el certificado de aptitud ambiental, «al no haberse realizado la presentación de dicha documentación ante la Municipalidad», y advierten que las obras en curso requieren permiso según el Código de Ordenamiento Territorial. El fiduciario de Ayres de Quitilipi responde el 6 de noviembre con su documentación: matrícula 5831, plano 1078, certificado de no inundabilidad (expediente 0090034-248087/2014) y certificado de aptitud ambiental (expediente 0090227-310880/2018).
\end{ficha}

\begin{ficha}{Expediente 0090227-229822/2022-0 de la Secretaría de Ambiente, 19 de octubre de 2022}
\textbf{Registro municipal de loteos} (nota 83/22 del subsecretario Rodrigo Matías Páez al Fiscal de Estado, 19/09/2022), sobre el área de la microcuenca:
\begin{itemize}[leftmargin=*,nosep]
\item El Nogalar, etapas I y II: \textbf{irregular}; matrícula 6557; radicado en 2012.
\item La Milagrosa: \textbf{irregular}; matrícula 5800.
\item El Durazno, etapas I, III y IV: aprobadas en 2022; radicado en 2009. (La Resolución 646D/21 había aprobado las etapas 1, 2 y 4: el capítulo \ref{cap:loteo} registra la discordancia.)
\item El Durazno, etapas II y V: \textbf{irregulares}; matrícula 6558; 2019 y 2018.
\item Club de Campo Potreros de La Caldera: \textbf{irregular}; matrícula 4207; 2013.
\item El Nogalar IV: \textbf{irregular}; «zona aluvional»; 2012.
\end{itemize}

\textbf{Inspección conjunta del 12 de octubre de 2022} (Ambiente, con el Ing.\ Guillermo Fuchs por Recursos Hídricos y el subsecretario municipal), sobre la margen derecha, entre el acceso a Potreros de La Caldera y la bomba de agua de El Durazno. Constató \textbf{dentro de la línea de ribera}: un alambrado paralelo a la Avenida Costanera, viviendas en lotes cerrados, pilares eléctricos, defensas precarias de ramas y áridos, la propia avenida y sus calles, «mallas para tamiz de áridos, que evidencia actividad de explotación minera» y la bomba de agua de El Durazno. Una retropala acordonaba material forestal en el cauce; el operario dijo que hacía defensas para El Durazno, que el material venía del cambio de uso de suelo de su etapa 3, y no exhibió autorización. Dentro del loteo: calles recién abiertas con volteo de vegetación, postes de alumbrado, zanjas para la red de agua y viviendas. La comparación de imágenes de 2018 y de agosto de 2022 muestra un «claro proceso de urbanización con la consecuente reducción de la cobertura boscosa».

\textbf{Plan de fiscalización:} intimar, por medio del municipio, a los loteos informados y al loteo abierto sobre las matrículas 1555 y 161 a presentar en treinta días una auditoría ambiental; sumario administrativo ya iniciado por este último.

\textbf{Denuncia penal SCLC-283/2022} (19 de noviembre de 2022): foco de incendio en el catastro 1555, sobre la Ruta 9, inmueble con sumario abierto por desmonte ilegal.
\end{ficha}

\begin{cautela}
\textbf{Es el primer documento oficial del expediente que califica de irregulares, con matrícula, los loteos que este libro sigue}, y lo firma el propio municipio. Dos años antes, en 2020, ese municipio ya había intimado a varios emprendimientos porque no le habían presentado la documentación.

Y hay una coincidencia de fechas: \textbf{las etapas III y IV de El Durazno figuran aprobadas en 2022}, el mismo año en que la inspección encontró desmonte en la etapa 3 y su material acordonado en el cauce, dentro de la línea de ribera \pendiente{sumario y auditorías}.
\end{cautela}

\subsection*{El financiamiento del estudio}

\begin{ficha}{Expediente 0090034-224086/2022-0, octubre de 2022}
El 12 de octubre de 2022, el secretario de Recursos Hídricos pide al representante de la Provincia ante el Consejo Federal de Inversiones, Ing.\ Ricardo Guillermo Villada, que gestione fondos del CFI para estudios en la cuenca del río La Caldera, «a fin de dar cumplimiento a la sentencia». El 14 de octubre, Villada autoriza el proyecto, a iniciarse una vez presentados el proyecto de ejecución y la documentación de la consultora.
\end{ficha}

Si el contrato con la consultora que elaboró el Plan se pagó con esos fondos, los documentos disponibles no lo dicen \pendiente{fuente de financiamiento}.

\subsection*{La presentación del estudio}

\begin{ficha}{«Plan Maestro de Drenaje Pluvial--Fluvial con Proyecto Licitatorio --- La Caldera». Descripción general e inicial del estudio}
Veintidós diapositivas, sin fecha, con los logos del Ministerio de Producción y Desarrollo Sustentable y de la Municipalidad de La Caldera. Objetivo: «contribuir de manera concreta en la disminución de la vulnerabilidad de la localidad de La Caldera y sus áreas de influencia ante las inundaciones provocadas por eventos pluviales».

\textbf{Diagnóstico:} «en los últimos 20 años, la localidad de la Caldera ha tenido una \textbf{expansión urbana hacia el sur} \ldots\ no acompañada con el desarrollo de infraestructura vial, pluvial y de defensa fluvial necesaria»; «áreas que anteriormente eran predios productivos fueron desarrollados para emprendimientos urbanos»; las quebradas que cruzan la trama urbana «no cuentan con obras hidráulicas de paso ni contenciones marginales»; y «la extracción de áridos y colocación de descartes de manera no planificada puede generar la desviación de los cursos hacia sectores vulnerables».

\textbf{Método anunciado:} siete subcuencas (Volcán--Yacones, Santa Rufina, San Alejo, Las Nieves, Lesser, Castellanos y Angosto; lámina \ref{lam:pmdp}), modelo digital de terreno, modelación hidrodinámica, un índice de peligrosidad hídrica que pondera impermeabilidad y pendiente (20\% cada una) y estado de los sistemas pluvial y fluvial (30\% cada uno), y un proyecto licitatorio para una obra priorizada.
\end{ficha}

\begin{cautela}
\textbf{El diagnóstico coincide con el de este libro, y lo dice el propio Estado.} Expansión hacia el sur, predios productivos convertidos en loteos, quebradas sin obras y áridos que desvían el cauce: son los cuatro hechos que los capítulos \ref{cap:loteo}, \ref{cap:ribera} y \ref{cap:resistencias} reconstruyen por separado.

\textbf{Pero varias diapositivas no muestran La Caldera.} Las de «estudios base» y «modelación hidrodinámica» presentan un poblado en paisaje árido, y el modelo de terreno de la diapositiva 15 va de \textbf{2.916 a 3.023 metros} (lámina \ref{lam:pmdpcotas}). La Caldera es un valle de yungas a unos 1.400 metros. Lo más probable es que sean ejemplos de otro trabajo, usados para mostrar el método; la presentación no lo aclara, y el título interno del archivo ---«Servicios de Ingeniería para Contratistas de PPP Viales»--- sugiere una plantilla reutilizada. Una presentación de anuncio puede ilustrar con otros casos. Lo que corresponde verificar es que el Plan final, el de doscientas treinta y cuatro hojas, sí modele la cuenca de La Caldera.
\end{cautela}

\lamina{lam-pmdp-subcuencas.jpg}{lam:pmdp}%
{Las siete subcuencas del Plan Maestro de Drenaje}%
{\textbf{Diapositiva 4 de la presentación del Plan Maestro de Drenaje Pluvial--Fluvial.} Sobre imagen satelital, las subcuencas Volcán--Yacones, Santa Rufina, San Alejo, Las Nieves, Lesser, Castellanos y Angosto, con la red de drenaje en celeste. La del Angosto, la mayor, desagua hacia el este, en la dirección del Mojotoro.}%
{Presentación «Plan Maestro de Drenaje Pluvial--Fluvial con Proyecto Licitatorio --- La Caldera. Descripción general e inicial del estudio», Ministerio de Producción y Desarrollo Sustentable y Municipalidad de La Caldera, sin fecha. Diapositiva completa, sin alteraciones.}

\lamina{lam-pmdp-cotas.jpg}{lam:pmdpcotas}%
{Un modelo de terreno de 2.916 a 3.023 metros}%
{\textbf{Diapositiva 15 de la misma presentación, bajo el título de «estudios base».} A la izquierda, un modelo digital de terreno cuya escala va de 2.916,2 a 3.023,0 metros sobre un poblado en paisaje árido; a la derecha, la red de calles de ese mismo poblado. Ni la altitud ni el paisaje corresponden a La Caldera.}%
{Misma presentación, diapositiva 15. Se recortó el margen blanco; no se alteró la imagen.}

\section{La causa sigue abierta}

Este libro dio por cerrado el litigio con la aprobación del Plan de Manejo en febrero de
2026. \textbf{No estaba cerrado.} La consulta del sistema de expediente digital del Poder
Judicial de Salta, en septiembre de 2026, muestra que la ejecución continúa y permite
reconstruir la cronología con precisión de día.

\begin{ficha}{Expediente MIN 780280/22 (17-00780280-0) --- Juzgado de Minas, Distrito Centro}
Carátula: \textit{Lazarte Vigabriel, Eduardo José Ignacio contra Municipalidad de La Caldera;
Provincia de Salta por otros}. Jueza: \textbf{Victoria Mosmann}. Iniciado el \textbf{26 de julio de
2022}; Secretaría Nº 2; tribunal de segunda instancia, la Corte de Justicia.

\textbf{Sujetos registrados:} el actor, Eduardo José Ignacio Lazarte Vigabriel; la Provincia de
Salta, con la Dra.\ Paula Gabriela Herrando; la Municipalidad de La Caldera, con su apoderado,
Héctor Alfredo Velásquez; la Fiscalía Civil, Comercial y Laboral Nº 2, Dr.\ Agustín Vidal; y
\textbf{tres dependencias de la Secretaría de Minería} ---Catastros, Mesa de Entradas y Dirección
General de Minería---. \textbf{Jardín Celestial S.R.L.\ no figura entre los sujetos de este
expediente}: su carácter de tercero consta en la carátula del incidente INC 23860/1 y en la de
la contratación de la consultora, no en la ejecución.

\textbf{Nov.\ 2022} --- se acompaña la \textbf{Línea de Base Socio Ambiental}.
\textbf{Sept.\ 2023} --- se acreditan \textbf{constancias de consulta libre, previa e
informada}.
\textbf{Nov.\ 2023} --- se acredita en tres presentaciones el \textbf{taller participativo} del 10 de octubre.
\textbf{Dic.\ 2023} --- informe final y documentación, puestos a disposición de las
\textbf{veedoras}.
\textbf{Feb.\ 2024} --- \textbf{denuncia de incumplimiento}; se acusa \textbf{rebeldía}.
\textbf{6 y 10 de dic.\ 2024} --- se acompaña el \textbf{Plan de Manejo}, en digital y en
papel.
\textbf{30 de dic.\ 2025} --- llamado de autos para interlocutorio, con habilitación de feria.
\textbf{5 de ene.\ 2026} --- \textbf{denuncia de incumplimiento} y pedido de audiencia
urgente, cargada por el portal a las 10:48 y proveída el 6.
\textbf{7 a 14 de ene.\ 2026} --- se acompañan informes, planillas, partes del Servicio Meteorológico Nacional y partes diarios del 1 de diciembre al 8 de enero; vista fiscal; un escrito que contesta la vista e impugna; y, el 14, un nuevo informe del actor que «advierte sobre» un decreto del juzgado.
\textbf{19 de ene.\ 2026} --- \textbf{denuncia de los hechos del 18 y 19 de enero} y pedido de
\textbf{declaración judicial de emergencia en la zona}.
\textbf{12 de feb.\ 2026} --- \textbf{resolución interlocutoria que aprueba el Plan de Manejo},
notificada al día siguiente.
\textbf{25 de mar.\ 2026} --- se oficia al \textbf{Registro de Procesos Colectivos}.
\textbf{Mayo--jun.\ 2026} --- pedido y concesión de \textbf{prórroga}; se tiene por cumplida.
\textbf{4 de sept.\ 2026} --- \textbf{última actuación registrada}: cédulas de notificación.
\end{ficha}

Cinco cosas cambian con esto, y dos son correcciones a lo que este libro afirmaba.

\textbf{Primera, y es una corrección: el Plan tardó catorce meses en aprobarse desde que se
presentó.} Se acompañó el 6 de diciembre de 2024 y se aprobó el 12 de febrero de 2026. Este
libro atribuía la demora al proceso de elaboración; buena parte de ella transcurrió con el
Plan ya presentado y en el juzgado.

\textbf{Segunda, y también es una corrección: hubo consulta previa.} En septiembre de 2023 se
acreditaron «constancias de consulta libre, previa e informada», y en noviembre las constancias del taller
participativo del 10 de octubre. Este libro declaraba, en una versión anterior del capítulo \ref{cap:metodo}, que no pudo verificarse la existencia de comunidades ni de consulta; ese capítulo ya lo corrigió para las comunidades, y esta pieza lo corrige para la consulta. \textbf{La consulta existió y consta en el
expediente}, y el acta del 12 de septiembre, transcripta más arriba, dice a quiénes se consultó y qué plantearon.

\textbf{Tercera: hubo veedores judiciales.} El expediente registra informes de veedores y
documentación puesta a su disposición: son las veedoras del COPAIPA cuyas observaciones al diseño del taller este capítulo transcribe.

\textbf{Cuarta: se denunció incumplimiento dos veces}, en febrero de 2024 y el 5 de enero de 2026,
y en una de ellas se acusó rebeldía. La denuncia del 28 de diciembre que registra el expediente administrativo (véase más abajo) no es una tercera: en el judicial no hay presentación del actor entre el 26 de diciembre y el 4 de enero. La segunda es de cinco semanas antes de que se aprobara el Plan, y dos semanas más tarde llegó la denuncia por los hechos del 18 y 19 de enero.

\textbf{Y quinta: la causa sigue abierta.} La última actuación registrada es del 4 de
septiembre de 2026. Entre la aprobación del Plan y esa fecha hubo oficios, un pedido de
prórroga concedido, informes de cumplimiento y la inscripción en el Registro de Procesos
Colectivos.

\begin{cautela}
Lo que antecede se apoya en la \textbf{lista de actuaciones} del sistema de expediente
digital: títulos, fechas, números y firmantes. \textbf{No se leyó el contenido de esas actuaciones} ---salvo el Plan y el acta del 12 de febrero, que este capítulo analiza más arriba---, ni los informes de los veedores.

De modo que este libro puede afirmar \textbf{qué se hizo y cuándo}, y no \textbf{qué dice cada
pieza}. Los informes de los veedores y las presentaciones posteriores a febrero de 2026 figuran en el apéndice \ref{cap:pedidos}, en la sección del Juzgado de Minas.
\end{cautela}

Y hay una última observación, que es sobre este libro y no sobre la causa. La información
anterior estuvo siempre disponible: el sistema de expediente digital es público y la consulta
lleva minutos. \textbf{Este trabajo reconstruyó durante meses, por prensa y por actos
administrativos, una cronología que el propio expediente publica en una tabla.} Queda
consignado entre los errores corregidos de esta investigación.

\section{Enero de 2026: el ciclo se repite}

Siete años después de los aluviones que originaron el amparo, y semanas antes de la aprobación judicial del Plan, las crecidas volvieron. \textit{Que Pasa Salta} consignó que, pasando el loteo El Durazno, el río se desbordó sobre el camino y lo socavó; que la zona más afectada va desde el arroyo Guaranguay hacia el sur, dejando aisladas a decenas de familias; que hubo casas inundadas y pérdidas materiales; y que los vecinos venían reclamando durante meses, con un tapado provisorio que no resistió.

\textit{El Tribuno} documentó una obra provincial: un badén construido por debajo del lecho del arroyo Guaranguay --- cauce seco casi todo el año pero que al crecer aísla los barrios del sur --- por ciento veinte millones de pesos, con cartel de obra provincial, expediente 366-16303/25, plazo de sesenta días y la Municipalidad de La Caldera como contratista. Los vecinos sostuvieron que el badén no era la solución y que correspondía encauzar el río y construir gaviones.

\subsection*{Lo que dicen los expedientes de la emergencia}

\begin{ficha}{Expediente 0090034-115343/2020 (Fiscalía de Estado y Recursos Hídricos), expediente 0110366-1348/2026 (Secretaría del Interior) y documentos municipales, diciembre de 2025 a enero de 2026}
\textbf{Diciembre de 2025.} Los partes diarios de Defensa Civil repiten alertas amarillas por tormentas para la «zona baja» y los «valles» de La Caldera.

\textbf{28 de diciembre.} Un temporal provoca una \textbf{crecida violenta del arroyo Guaranguay}; según \textit{El Tribuno}, los vecinos evitaron que arrastrara una camioneta; \textit{Opinorte}, al día siguiente, habla en cambio de un vehículo arrastrado (véase más abajo). Las dos notas no coinciden, y este libro no encontró el hecho en un documento oficial. El amparista denuncia un hecho nuevo y el \textbf{incumplimiento de la medida cautelar}: no hubo medidas preventivas pese al alerta del Servicio Meteorológico Nacional, los bomberos no tenían equipamiento y el municipio se limitó a avisar en redes, con la crecida en curso, que no se podía pasar. Pide habilitar la feria judicial. \textit{En el expediente judicial ese escrito no figura con esa fecha: la lista de actuaciones no registra presentaciones del actor entre el 26 de diciembre y el 4 de enero; el 30 de diciembre el juzgado dictó el llamado de autos con habilitación de feria sin que la lista registre un pedido previo del actor, y la denuncia de incumplimiento entró por el portal el lunes 5 de enero.}

\textbf{29 de diciembre.} Recursos Hídricos recorre la zona y elabora un legajo de \textbf{«Obra de emergencia: Río La Caldera»}, fechado en diciembre, con un \textbf{presupuesto oficial de \$27.831.927,50}: un terraplén de 180 metros en el barrio Sabaleta ---así en el legajo hídrico; los informes municipales escriben \textbf{Zavaleta} y el segundo informe de Recursos Hídricos, \textbf{Zabaleta}: es el mismo barrio y este libro no puede establecer cuál es la grafía correcta---, el encauzamiento del río en dos tramos de 160 metros frente a Santiago Apóstol ---arrastrando con topadoras material del tercio medio---, y un terraplén contra desbordes de 160 metros en El Nogalar 4.

\textbf{3 y 4 de enero de 2026.} Lluvias intensas inundan Nogalar 4, El Durazno y Potrero de La Caldera. Esa madrugada, en el grupo del Comité Operativo de Emergencia, el intendente informa que la principal zona afectada es «Nogalar 4, loteo Potrero y ruta nacional 9 a la altura de km 1618», que no hay riesgo de vida y que la máquina municipal debe mover material para poder llegar. El 5, el área de Desarrollo Social llena fichas de relevamiento de emergencia en Nogalar 4.

\textbf{5 de enero.} La Municipalidad pide a la Secretaría del Interior \textbf{\$10.000.000} para encauzar el río en la zona sur, y a la Secretaría de Obras Públicas, asistencia financiera urgente para una topadora, una retroexcavadora y un camión volquete: «las zonas más comprometidas son Nogalar 4, Potrero y Wierna». Ese día entran al sector una retroexcavadora y una topadora de Vialidad de Salta, puestas al servicio del municipio. El 7, el municipio pide además al Ministerio del Interior de la Nación que priorice su solicitud de maquinaria en el programa «Municipios de Pie».

\textbf{6 de enero.} La jueza de feria, María Guadalupe Villagrán, \textbf{intima a la Municipalidad y a la Provincia a informar en un día} las medidas adoptadas; pregunta si el municipio activó el Protocolo para Crecidas e Inundaciones que obra en el expediente y si Defensa Civil trabaja en la zona; y \textbf{ordena cumplir el protocolo}, con intervención de Defensa Civil si la capacidad municipal se supera, «bajo apercibimiento de incurrir en desobediencia judicial».

\textbf{7 de enero.} Fiscalía de Estado remite la intimación a Recursos Hídricos, a la que todavía ubica en el Ministerio de Producción y Minería. La respuesta llega el mismo día, firmada por el coordinador de la \textbf{Subsecretaría de Recursos Hídricos, Jefatura de Gabinete de Ministros}: acompaña el legajo de emergencia y describe como «medida concreta adoptada y en ejecución» la asistencia técnica en el lugar.
\end{ficha}

\begin{cautela}
\textbf{Las cifras de la emergencia, puestas al lado de las del Plan.} Para el tramo urbano del río, la Provincia presupuestó en diciembre de 2025 una obra de emergencia de 28 millones de pesos, con terraplenes de material del propio cauce. En el mismo expediente judicial, el Plan de Manejo que se aprobaría semanas después cifra la defensa definitiva del río La Caldera en 13.697 millones, a precios de 2024. El municipio, por su parte, pedía 10 millones para máquinas.

\textbf{Y la intimación judicial del 6 de enero pregunta lo mismo que este libro}: si el protocolo existía para usarse. Las denuncias del 5 y del 19 de enero que registra el expediente son la respuesta de los amparistas.
\end{cautela}

\subsection*{Lo que informó cada uno después}

\begin{ficha}{Documentos municipales, provinciales y registrales, enero de 2026}
\textbf{Comunicación oficial del municipio (4 y 5 de enero).} El barrio Nogalar 4 fue el más afectado; de treinta domicilios visitados, «sólo una familia tuvo pérdidas». Y el intendente declara que los anegamientos en El Nogalar y otros barrios del sur son \textbf{«consecuencia directa de habilitaciones de loteos hechas entre diez y veinte años atrás, que permitieron urbanizar zonas en forma irregular, que no estaban preparadas, que están dentro del río, y sin exigir estudios de no inundabilidad, ni impacto ambiental ni infraestructura esencial»}; que los pedidos de defensas «reciben negativas, porque los organismos no aprueban hacer defensas dentro del río»; y que el municipio reiterará a los loteadores que deben cumplir los trámites.

\textbf{Informe de la Dirección municipal de Desarrollo Social (28 de enero).} El 3 de enero a las 20, el intendente declara el alerta por lluvias en la zona alta (Santa Rufina y San Alejo); a las 22 llega la crecida; a las 23 la policía informa que no puede llegar a El Nogalar IV por la crecida de los arroyos. Zonas afectadas: El Nogalar IV y el club de campo Potreros de La Caldera; la Ruta 9, con derrumbes en el km 1618. El hogar de ancianos se prepara para recibir evacuados. El río, desviado sobre su margen, \textbf{socavó e inundó unos 600 metros de la avenida Costanera} y dejó sin ingreso ni egreso a El Nogalar y a Potreros. El agua entró \textbf{en los terrenos de catorce viviendas} ---otras seis se salvaron por sus muros perimetrales y el relleno de sus lotes--- y se asistió a \textbf{ocho familias}. En Potreros de La Caldera, los asistidos declararon no vivir allí: habían alquilado la casa por el fin de semana. La Provincia envió quince colchones, veinte frazadas, treinta bidones de agua, treinta módulos alimentarios y dos rollos de plástico.

\textbf{Informe municipal de obras (3 al 27 de enero).} La noche del 3, la retropala municipal y una pala cargadora de un particular llegaron al lugar, pero «no se pudo actuar debido a que la maquinaria no es la adecuada»; la maquinaria pesada llegó el 5 al mediodía. Intervinieron Vialidad Provincial, Recursos Hídricos, Obras Públicas, la empresa Valle Vial S.R.L.\ y el municipio: una topadora (180 horas), dos retroexcavadoras de oruga (100 horas cada una), otra de Valle Vial (50), la retropala municipal (70) y camiones (65 cada uno). Se hicieron bateas y cordones con material del propio río en unos 1.000 metros en Nogalar 4, más 150 metros de espigones; 550 metros de bateas en el camping y el barrio Zavaleta; 200 metros de defensa con un terraplén de 80 en Santiago Apóstol; un terraplén de 100 metros en la Costanera, y 120 metros de espigones de enrocado junto al cementerio. Para el futuro, el municipio proyecta espigones de material del río que «luego serán de neumáticos en desuso». Una planilla anexa declara \textbf{\$8.311.560,09 de combustible} usado sólo por las máquinas entre el 5 y el 27 de enero; la suma cierra, pero el mismo número de comprobante (0035-00006919) figura en dos fechas distintas y la numeración salta de la fila 19 a la 22.\pendiente{comprobantes de combustible de enero de 2026, en particular el 0035-00006919, que figura en dos fechas, y filas 20 y 21 de la planilla municipal}

\textbf{Informe municipal de gestiones (fechado por error el «27 de enero de 2025»).} Enumera lo que el municipio pidió sin obtener respuesta o autorización expresa: un informe de impacto ambiental para intervenir el río en su área de exclusión (expediente 90302-250447/2023), sin observación alguna al 26 de enero de 2026; una nota a Recursos Hídricos de noviembre de 2023, respondida en diciembre como favorable «sin que ello implicara autorización expresa»; un proyecto de defensas para Santiago Apóstol y Nogalar 4, presentado al Ministerio de Infraestructura el 21 de febrero de 2024 (expediente 386-219264); una nota a Minería del 10 de julio de 2025 (registrada el 21) para encauzar el río junto al puente de acceso ---«en el verano pasado el río se llevó parte del paseo peatonal de la costanera»--- y «en cualquiera de sus tramos», acompañada de los títulos de la maquinaria municipal (un camión de 2002 y una retroexcavadora de 2011), sin respuesta; y dos pedidos a la Secretaría del Interior, en julio de 2025 y en enero de 2026, de fondos para \textbf{pagar los honorarios judiciales del amparo}, que el municipio dice no poder solventar.

\textbf{Segundo informe de obras de Recursos Hídricos (28 de enero).} En el barrio Zabaleta, «este año colapsó el caño y la cámara del acueducto norte», que Aguas del Norte ya reparó. El 9 de enero se sumaron Obras Públicas, Vialidad, la Municipalidad de la ciudad de Salta y empresas particulares. El 18 hubo un nuevo evento y los sedimentos anegaron los badenes. El 21, equipos de Aguas del Norte y uno particular, a pedido de Obras Públicas, encauzaron el río aguas arriba del puente, \textbf{«despejando la margen derecha donde se vio afectado el acueducto norte»}. El 23, la verificación se hizo con un vuelo de dron, «equipo particular perteneciente al personal» de la Subsecretaría. El 27 se terminaron ese encauzamiento y la reparación del acueducto, y comenzaron las defensas de Santiago Apóstol.

\textbf{Dirección General de Inmuebles (22 de enero).} Informa al Juzgado de Minas la \textbf{inscripción definitiva del embargo sobre la matrícula 1.846}, de titularidad de la Municipalidad, ordenado en la causa «Ejecución de honorarios; Chávez, Julio; Lazarte Vigabriel\ldots\ contra Municipalidad de La Caldera» (expediente 938.744/25).
\end{ficha}

\begin{cautela}
\textbf{Cinco lecturas.} La primera: el diagnóstico que este libro reconstruye con planos y expedientes lo formula ahora el propio intendente, con las mismas palabras ---loteos habilitados hace diez o veinte años, dentro del río, sin estudios---. El registro municipal de 2022 ya los calificaba de irregulares.

La segunda: la misma comunicación que lo reconoce reduce el daño a «una familia», mientras el informe interno del municipio cuenta catorce terrenos inundados, ocho familias asistidas y dos barrios aislados.

La tercera: el municipio pidió en 2023 y en 2025 permiso para intervenir el río con su maquinaria y no lo obtuvo; en 2023, Minería lo había paralizado por hacerlo sin permiso. En enero de 2026, en emergencia, el río fue intervenido por cinco organismos y dos empresas, con material del propio cauce.

La cuarta: en el tramo del puente, el encauzamiento se orientó a despejar la margen donde se había dañado el acueducto que abastece a la capital, y las defensas de Santiago Apóstol empezaron cuando esa reparación terminó.

La quinta: \textbf{el inmueble que la Provincia donó en 2011 «exclusivamente al desarrollo integral de la zona» (Ley 7699) está embargado} para pagar los honorarios del amparo sobre el río. Al día de hoy, el municipio no pudo pagarlos y la Provincia no respondió su pedido de fondos \pendiente{estado de la ejecución}.
\end{cautela}

\subsection*{Lo que el expediente agrega de 2023 a 2026}

\begin{ficha}{Transcripción integral del legajo del amparo y actuaciones conexas (75 documentos, 2022--2026)}
\textbf{16 de agosto de 2023.} La Dirección de Infraestructura Hídrica inspecciona, a pedido del municipio, la margen derecha junto al camping municipal. Encuentra un terraplén \textbf{construido por la cantera La Serena}, que sirve de acceso a sus camiones y, a la vez, de defensa del camping. El municipio quiere prolongarlo hacia el sur. La Provincia \textbf{recomienda no cortar el monte bajo}, más alto que el río y defensa natural del camping, y advierte que, por tratarse de una curva, la obra dirigiría el cauce contra la margen izquierda, donde ya hay erosión y reparaciones de la barranca de la Ruta 9.

\textbf{5 de diciembre de 2023.} Ante el pedido municipal de limpiar y encauzar ríos y arroyos, Recursos Hídricos responde que aplica el Código de Aguas y que \textbf{no le corresponde ser «ejecutor de proyectos y obras hidráulicas»}; que, según el Decreto 2299/03, las aguas pluviales que discurren por lugares públicos son responsabilidad municipal; y que considera favorable la tarea. Al día siguiente, el intendente pide a Desarrollo Social tanques, camas, colchones y alimentos para las familias de las zonas críticas ante el verano.

\textbf{12 de marzo de 2024.} Acta del Comité de Emergencia en la Municipalidad: el intendente explica el protocolo; se resuelve modificarlo y se entrega \textbf{un nuevo protocolo a los concejales para su aprobación por ordenanza} y su elevación a Defensa Civil; un concejal queda a cargo de activarlo. Los bomberos voluntarios informan que \textbf{«no cuentan con movilidad»} y que, ante una emergencia, acuden a la policía y al municipio. Se ofrecen lugares para albergar a noventa personas.

\textbf{Julio de 2025 y 5 de enero de 2026.} El intendente pide a la Secretaría del Interior, y luego reitera, fondos para pagar los honorarios regulados por la jueza de Minas: \textbf{\$10.004.770}, en la causa «Lazarte Vigabriel c/ Municipalidad de La Caldera y Provincia de Salta; tercero: Jardín Celestial S.R.L.». Las costas las impuso la Corte de Justicia, \textbf{en forma solidaria a la Provincia y al municipio}. La reiteración menciona el embargo preventivo sobre bienes municipales y que Fiscalía de Estado también pidió esos fondos.

\textbf{Enero de 2026.} La Municipalidad resume su gasto en la emergencia: \textbf{\$8.311.560,09 en combustible, \$11.890.000 en horas máquina y \$1.332.146 en personal, \$21.533.706,09 en total}. El respaldo incluye remitos de la estación de servicio ---que llevan, como es habitual, la leyenda «documento no válido como factura»---, facturas electrónicas de un proveedor local por el alquiler de una retroexcavadora (64 horas a \$121.000 más el traslado: \$7.986.000), recibos de tesorería y planillas de horas extra de operarios de Vialidad. El remito del 21 de enero coincide en número con la planilla de combustible, pero cobra el gasoil a \$1.808 y no a \$1.806, e incluye cuarenta litros de nafta que la planilla, limitada a las máquinas, no computa.
\end{ficha}

\subsection*{Lo que dijo el intendente}

\begin{ficha}{Declaraciones públicas del intendente de La Caldera, julio de 2025 a enero de 2026, compiladas en el legajo}
\textbf{3 de julio de 2025.} Ante la denuncia del diputado Luis Mendaña en la Cámara de que los intendentes del departamento «comenzaron a devolver obras», el intendente y el ministro de Infraestructura, Sergio Camacho, responden que un convenio de defensas firmado en 2023 por \$63 millones costaba ya más de \$200 millones, y que se lo rescindió para firmar uno nuevo con valores actualizados.

\textit{El índice de resoluciones del Boletín permite contrastar la denuncia y la respuesta. Las rescisiones de obras con la Municipalidad de La Caldera son de marzo de 2024: la defensa de Santiago Apóstol y Nogalar~4 y la ampliación de la red de agua del mismo barrio (Resoluciones 314/24 y 315/24, capítulo \ref{cap:defensas}). En 2025 el Boletín no publica ninguna rescisión con La Caldera; las que publica son de Vaqueros: los taludes del SUM de San Nicolás (247/25, mayo) y, el 24 de julio ---tres semanas después de la sesión---, la optimización de redes y el pavimento de la calle Cortázar (367/25 y 368/25), de las que en septiembre se aprueban legajos nuevos (482/25 y 483/25). La cifra del ministro, «más de 200 millones», es la del convenio nuevo: la Resolución 951/24, de octubre de 2024, rehízo la misma obra ---ciento ochenta metros de gaviones en Santiago Apóstol y El Nogalar~4--- por \$216.819.417,19, y la Resolución 479/26 la rescindió en septiembre de 2026 porque el cauce había cambiado (capítulo \ref{cap:defensas}).}

\textbf{29 de diciembre de 2025} (\textit{Opinorte}, sobre una entrevista radial). Tras el vehículo arrastrado en el badén del Guaranguay: «Hay loteos que se habilitaron hace 10, 15 o 20 años sin infraestructura, sin estudios y sin exigir nada a los loteadores». Atribuye algunos de esos emprendimientos a gestiones anteriores ---«Hubo sociedades, convenios y donaciones de tierras. Se usó maquinaria municipal, se abrieron calles»---, dice haber creado al asumir un registro de loteadores con exigencia de certificados de no inundabilidad, y pregunta por qué quienes los habilitaron y luego fueron legisladores no gestionaron las obras.

\textbf{5 de enero de 2026} (\textit{El Tribuno}). En El Nogalar~4, el loteo Potrero de La Caldera y Wierna viven unas treinta familias; a cinco les entró agua en el patio y una tuvo pérdidas. La defensa de ese sector figuró en planes y presupuestos, pero «el presupuesto presupone» y no hubo recursos. Consultado sobre El Nogalar ---que el periodista describe como loteos municipales---, describe el mecanismo: \textbf{«Se habilita el loteo privado, se donan parte del remanente sobre el río a la Municipalidad. La Municipalidad, por la necesidad de la gente, otorga esos terrenos. Luego, la donación de calle del loteo privado la hace el Municipio, que la acepta. Y después el Municipio tiene la obligación de hacer el acondicionamiento»}. Lo atribuye a las gestiones de Calabró, Mendaña y Escalera, y concluye: \textbf{«el Estado está financiando loteos privados»}. Sobre los barrios nuevos, dice que muchos lotes no pagan tributos porque «todavía están en vía de legalización».

\textbf{6 de enero} (\textit{Informate Salta}). La nota consigna que La Caldera tiene barrios autorizados en sectores que, según dictámenes técnicos de Recursos Hídricos, están en el lecho del río.

\textbf{19 de enero} (\textit{Que Pasa Salta}). Tras los rescates informados por la Policía: «Gracias a Dios, el agua no ingresó al domicilio de ningún vecino».

\textbf{21 de enero} (\textit{Informate Salta}). En Casa de Gobierno propone represas de retardo en la cuenca alta, más control de la extracción de áridos, más exigencias a los loteos y restricciones a los asentamientos en el lecho del río; el jefe de Gabinete coincide y sobrevuela la zona.
\end{ficha}

\begin{cautela}
\textbf{El mecanismo que este libro reconstruye con planos y matrículas lo describe ahora el intendente en una entrevista}: el loteador cede al municipio el remanente sobre el río, el municipio lo adjudica a familias, acepta las calles y hereda la obligación de mantenerlas. La infraestructura que el emprendimiento no hizo la paga el Estado, y los lotes, mientras no se regularizan, tampoco tributan.

\textbf{Las atribuciones a gestiones anteriores son afirmaciones de una parte en una disputa política}, formuladas por el intendente en funciones sobre sus antecesores ---uno de los cuales, como diputado, lo había acusado meses antes en la Legislatura---. El legajo no contiene la respuesta de los aludidos, y este libro las registra como declaraciones, no como hechos probados \pendiente{actos de habilitación}.

\textbf{Y las cifras del daño cambian según quién habla.} El intendente dijo que a una sola familia le entró el agua con pérdidas (5 de enero) y que no entró en ninguna vivienda (19 de enero). El informe de su propia Dirección de Desarrollo Social registra catorce terrenos con agua, ocho familias asistidas y dos barrios aislados, y la Policía informó rescates.
\end{cautela}

\section{La reunión de Casa de Gobierno}

El 21 de enero de 2026, en plena emergencia, el intendente participó en Casa de Gobierno de una reunión con el Gobierno provincial, los intendentes del Valle de Lerma y legisladores. Según el comunicado del propio municipio, el encuentro fue propuesto por el intendente y convocado por el jefe de Gabinete Sergio Camacho, con participación del ministro de Gobierno Ignacio Jarsún y de los intendentes de Vaqueros, Cerrillos, Rosario de Lerma, La Merced y Campo Quijano.

El comunicado consigna que el intendente ``propuso desarrollar un plan integral de manejo de la cuenca hídrica de La Caldera que solucione este problema a largo plazo'', señalando que las obras de defensa son paliativas. Y que el jefe de Gabinete se mostró de acuerdo con la creación de represas de retención, \textbf{pidió más control de extracción de áridos y que los municipios incrementen el control a los loteos}.

El Plan de Manejo Integral de la microcuenca no era una idea nueva en enero de 2026: era una condena judicial de diciembre de 2021 ---recurrida ante la Corte y seguida de una segunda sentencia en julio de 2023---, dictada contra el propio municipio y contra la Provincia, con criterio técnico expresamente definido. Tres semanas después de esta reunión, el mismo plan fue aprobado judicialmente.

Que el intendente proponga en 2026, como iniciativa propia y ante el Gobierno provincial, la elaboración de un plan al que una sentencia lo obligaba desde hacía cuatro años, es el registro más exacto de qué le ocurre a una condena de este tipo en la práctica administrativa. No se desobedece: se reabsorbe. Cuatro años después reaparece como propuesta de gestión, en un comunicado que no menciona la sentencia.

Sobre el criterio, una segunda observación. El tribunal ordenó priorizar infraestructura verde sobre obra civil. Lo que se propone y acuerda en la reunión son represas de retención: infraestructura gris. Puede ser técnicamente correcto y compatible con el plan aprobado, pero es la inversión del orden de prioridades que la sentencia fijó, y quien lo verifique deberá contrastarlo contra el texto del Plan.

Y el pedido del jefe de Gabinete a los municipios de incrementar el control a los loteos es la autoridad provincial reconociendo públicamente, en enero de 2026, el diagnóstico que el informe del PDUA había formulado en 2015 y que el Código reglamentó sin aplicarse.

\section{Febrero a mayo de 2026: el Plan aprobado y sin presupuesto}

\begin{ficha}{Acta de la audiencia del 12 de febrero de 2026 (Expte.\ MIN 780280/22) y expediente administrativo 0090034-115343/2020, febrero a mayo de 2026}
\textbf{La audiencia y lo resuelto} se detallan más arriba en este capítulo: aprobación del Plan, intimación al municipio bajo apercibimiento penal y cuarenta y cinco días corridos para que la Provincia presentara la planificación de las obras prioritarias.

\textbf{La Secretaría de Obras Públicas} (informe elevado el 20 de marzo) explica que las obras del Plan están desarrolladas \textbf{«a nivel de prefactibilidad»} y deben llevarse a proyecto ejecutivo. Actualiza su costo con el índice de la construcción del INDEC: de \$23.800 millones a \textbf{\$32.900.449.736,76 a enero de 2026}, de los cuales \$18.934 millones corresponden a la defensa del río La Caldera. Calcula que sólo los proyectos ejecutivos costarían el 4\,\% de ese monto, y que esos recursos \textbf{«no [están] contemplados en el presente ejercicio presupuestario»}. Describe la respuesta a la emergencia: retroexcavadoras dirigidas por Recursos Hídricos y una topadora de Vialidad que, entre el 5 de enero y el 28 de febrero, desplazó lateralmente \textbf{más de 30.000 m³} de material del cauce «que terminó por proteger íntegramente la margen derecha»; y un legajo de encauzamiento en cinco sectores, basado en un anteproyecto municipal, por \textbf{\$42.725.963,14} a valores de enero, a ejecutar en veinte días por ajuste alzado. Entre los objetivos de ese legajo figura que el encauzamiento «permite usar áreas cercanas al río para agricultura, infraestructura o urbanización». Y agrega dos constataciones: en la zona \textbf{«existen desarrollos que aún no cuentan con aprobación»} y que no diseñaron su drenaje, obras que deberían ejecutar los desarrolladores; y \textbf{«existen asentamientos, hoy con una gran densidad de construcción, que no respetaron la línea de ribera sancionada y se encuentran invadiendo el propio cauce del río»}. Una ortofoto adjunta marca junto a El Durazno una «zona sin planos aprobados aún». La Secretaría no asistió a la reunión convocada por Fiscalía para el 27 de febrero porque el expediente le llegó ese mismo día.

\textbf{27 de marzo.} Fiscalía presenta en el juzgado lo informado por Obras Públicas, y el expediente vuelve a circular para planificar el resto de las obras.

\textbf{1 y 15 de abril.} Obras Públicas ordena a Recursos Hídricos elaborar \textbf{«por administración»} ---con personal propio--- el proyecto ejecutivo, «priorizando habida cuenta de no estar contemplado en el Presupuesto 2026», de las \textbf{defensas del río La Caldera y de las defensas o encauces de los arroyos Guaranguay, Nogalar y El Durazno}, y coordinar con Minería el modo de explotación de áridos «debido a la sobreexplotación de áridos en el mencionado río, no regularmente llevados a cabo».

\textbf{28 de abril y 6 de mayo.} La Dirección General de Infraestructura Hídrica advierte que cumplir en plazo «devendrá en imposible», y Obras Públicas pide a Fiscalía de Estado \textbf{una prórroga de sesenta días}.
\end{ficha}

\begin{cautela}
\textbf{El Plan quedó aprobado y sin dinero para empezarlo.} Tres meses después de la aprobación, la Provincia no tenía presupuesto ni para los proyectos ejecutivos, que encargó a su propio personal, y pidió más plazo. La obra que sí se hizo fue la de emergencia: mover el material del cauce contra la margen derecha. Para dimensionarlo: según la planificación oficial difundida al aprobarse el Presupuesto 2026 ---en la misma sesión que la Ley 8.511---, la Provincia destina ese año \textbf{\$294.560 millones a obras}, 488 proyectos nuevos y la continuidad de más de 2.560, con un rubro genérico de «obras hídricas y de drenaje». El Plan de Manejo actualizado equivale a algo más del 11\,\% de esa cifra; sus proyectos ejecutivos, al 4\,\% del Plan, a unos \$1.300 millones, menos del medio por ciento. \textbf{Ninguno de los dos entró}, y la Ley 8518 de Presupuesto 2026 (B.O.\ Nº 22.092 del 15/12/2025) lo confirma. En sus planillas, el río La Caldera figura sólo en el «Plan de Trabajos Públicos 2026 --- Obras con financiamiento a confirmar», sin monto, con dos obras: «Defensas en Río Yacones zona de la Escuela» y la «construcción del camino alternativo sobre margen derecha del río La Caldera, tramo La Caldera hasta la intersección con la Ruta 9 en el puente Wierna». Con financiamiento asignado, en cambio, el departamento recibe la \textbf{reconstrucción del azud nivelador del dique Campo Alegre (\$541.501.257)} y \textbf{intervenciones de emergencia en su azud y su toma (\$191.143.863)} ---adjudicadas por la Resolución 568/25 de Obras Públicas, en noviembre de 2025, \textbf{a la Municipalidad de La Caldera}, por \$211.793.261,78 a valores de julio de 2025 y con cargo al ejercicio 2025 ---un 11\,\% más que la partida de 2026---; en mayo de 2026 la Resolución 285/26 le aprueba un adicional de \$52.477.584,20, casi un cuarto más, porque «las crecidas del caudal del rio generadas por las intensas lluvias» cambiaron «el volumen de material aluvional a remover en la zona de encauzamiento». La Provincia contrata así para mover áridos junto a la toma al mismo municipio al que en agosto de 2023 Minería había paralizado por extraerlos en el área de exclusión del acueducto---, ambas obras del sistema que abastece de agua a la capital, y la sección I de la Ruta 9 entre Salta y La Caldera (\$18.026.504.000, con aportes de Vialidad Nacional). Para toda la provincia, las «obras de defensa y encauzamientos en los diferentes municipios» suman \$451.251.048. \textbf{En el presupuesto, el agua de La Caldera que va a la ciudad tiene partida; el río que inunda sus barrios, no.} Y el mismo año, la Provincia sí consiguió financiamiento externo para proyectos ejecutivos en otro punto de la cuenca del Mojotoro: la Resolución 733D/26 del Ministerio de Economía (B.O.\ Nº 22.263 del 07/09/2026) adjudicó, con fondos del BIRF, los del Nodo Logístico y Parque Industrial de Güemes ---incluidos «hidrología, hidráulica y proyecto de desagües pluviales»--- por \textbf{USD 98.271}. En la lista corta de ese concurso figuraba \textbf{YSATÍ Ingeniería}, la consultora del Plan de Manejo, asociada a Aetos Consultores; no estuvo entre las cuatro propuestas que se abrieron.

\textbf{La observación que el párrafo anterior dejaba abierta puede ahora verificarse.} El Plan prioriza, en este orden, el camino de ribera, el bosque de ribera y la defensa del río La Caldera. Las dos primeras ---la primera, infraestructura de acceso; la segunda, la única medida verde del cuadro--- no figuran en lo que la Provincia mandó proyectar, que empieza por las defensas del río y de los arroyos. \textbf{El orden de la sentencia, que manda privilegiar la infraestructura verde, se invirtió en la ejecución.}

\textbf{Y el Estado dejó escrito lo que este libro documenta.} Obras Públicas reconoce desarrollos sin aprobación y asentamientos dentro del cauce, más allá de la línea de ribera sancionada; Recursos Hídricos admite que las obras de emergencia no son las del Plan; y el propio legajo de emergencia enuncia, entre sus fines, habilitar el uso urbano de las áreas cercanas al río.
\end{cautela}


